\chapter{Measuring School Culture: Warmth and Strictness in English Secondary Schools}
\label{ch:paper1}
\resetcounterformats

\begin{center}
  \textit{Damian Phelan%
    \thanks{Department of Economics, University College London.
      Primary supervisor: Michela Tincani. Secondary supervisor: Lars Nesheim.
      I thank Michela Tincani and Lars Nesheim, and seminar participants at
      the Institute for Fiscal Studies and UCL, for comments. This work was
      supported by the Institute for Fiscal Studies and the Stone Centre at
      UCL. The school visits and headteacher interviews were approved by
      the UCL Research Ethics Committee (project 22979.001).}}
\end{center}

\begin{chapterabstract}
How warm a school is towards its pupils, and how strictly it runs, are
widely thought to affect how much progress pupils make. This paper measures both, and asks
how much of each can be inferred from the documents that every school
produces. In \NTierOne{} schools a research team that I led collected a gold-standard measure: a day
of structured observation recorded what each school does, and an interview
with its headteacher recorded what the head says it does. I then carried out textual analysis of
every school's Ofsted inspection report, behaviour policy and website using two different approaches.
Firstly, a prediction model learns, from the visited schools, which words and phrases
are linked with warmer or stricter schools.  Secondly, a large language model (LLM) is asked to
use a one-page mark scheme to evaluate the document on each dimension on a 1 to 5 scale.
There are four findings. First, inspection reports carry a modest signal of how warm a
school actually is: the prediction model gives a correlation of $\PredWOfsted$. Second, we see a similar modest correlation
of $\RubricSOfstedDec$ for strictness measured by the LLM using the mark scheme.
Third, what a headteacher says about their school is not what observers find there, and the gap
is widest for warmth. The heads' ratings and the observers' scores correlate at $\GoldWarmthSplit$ for warmth and $\GoldStrictnessSplit$ for strictness. Correcting both for measurement error raises them to $\DisSplitW$ and $\DisSplitS$. Fourth, schools' own documents match what their heads say about strictness and carry no signal about warmth.
\end{chapterabstract}

% ---------------------------------------------------------------
\section{Introduction}
\label{sec:p1_intro}

How warm a school is towards its pupils, and how strictly it runs, are two of
the things most often said to affect how much progress its pupils make. In this paper I construct
measures of these two dimensions of school culture based on a structured day of observations at a school. I refer to these
two dimensions as warmth and strictness. I then ask how well each dimension can be predicted from publicly available texts that exist for all 3,332 secondary schools in the country.

The two dimensions have precise meanings here. Warmth is
the relational quality of a school: whether pupils are known, cared for and
supported by the adults around them. In a lesson it shows in whether the
teacher uses pupils' names, how praise is given and how the teacher and the
class speak to one another. Between lessons it shows in how staff and pupils
deal with each other in corridors and at break. It also shows in whether rewards are
used consistently across the lessons of the day. Strictness is the school's
behavioural structure: whether
expectations of conduct are clear, enforced and met. In a lesson it shows in
how much misbehaviour is tolerated before a sanction, how much disruption
occurs and how pupils respond when they are corrected. Outside lessons it
shows in how pupils arrive, how they move between lessons and how the dining
hall runs. Strictness in this sense means clear rules applied consistently rather than harsh punishment. A school whose
sanctions are applied by every member of staff in the same way is strict in
this sense.

The two dimensions come from the parenting typology of \citet{baumrind1971}
and \citet{maccoby1983}, which classifies parenting by demandingness and
responsiveness and finds that the combination of both, the authoritative
style, produces the best child outcomes across cognitive and non-cognitive
domains. \citet{majumder2015} and \citet{zhang2017} find the same pattern in educational attainment: children of authoritative parents do best.

\citet{pellerin2005} applies the typology to schools, classifying each school as authoritative,
authoritarian, permissive or indifferent based on its levels of demandingness
(academic and behavioural expectations) and responsiveness (support for pupils). She finds that
pupils in authoritative schools have lower dropout and truancy rates.
\citet{gregory2010,gregory2011} use a similar typology and find that more authoritative schools
have less bullying, less victimisation and lower suspension rates. \citet{gottfredson2005} finds
that a school's climate, the clarity and fairness of its rules and how consistently they are
enforced (demandingness), predicts the level of disorder in the school even after accounting for the demographics of its pupils.

All of these studies rely on surveys of pupils, parents or staff and
use slight variations on the demandingness and responsiveness language.
This paper refers to the dimensions as warmth and strictness, as this is common terminology in the
teaching profession, and measures them through direct observation.

What makes the two dimensions measurable in a school,
as opposed to in a family, is that both are visible in public. A researcher
sitting in a lesson can see whether names are used and how disruption is
dealt with, and a researcher standing in a corridor can see whether pupils
move calmly and whether staff speak to them. Neither requires access to
anything private. I treat warmth and strictness as two separate dimensions
rather than as the two ends of one scale, because a school can be warm
without being orderly and orderly without being warm. Whether warm schools
also tend to be strict, and strict schools warm, is an empirical question,
and scoring the two separately lets the data answer it.

I also observed a third dimension, the quality of teaching practice.
This is not part of culture as I define it, but it is the most obvious alternative channel
for any association between culture and results, and it is kept in the analysis for that reason.

Warmth and strictness are not directly observed for schools in England. However, there is public
text available for each school that may be informative about their culture.
Every school has a report written by an external inspector, a behaviour policy written by the
school, and a website written by the school. For a smaller set
of 103 schools this study adds data from direct observation of the school at work and from a
structured interview with its headteacher. For a further \NTierTwo{} schools we have the interview alone. Those schools were not visited, so no observer scored them, but they widen the sample for one later test: whether what a school writes about itself matches what its headteacher says.

From the structured observation and interview data, I construct a gold-standard measure of
each dimension. Because the interview and the visit disagree substantially about the same school,
the gold standard is in fact two measures rather than one: an \emph{enacted}
score, from what the observers recorded, and an \emph{espoused} score, from
what the headteacher reported.

I then ask whether these scores could be predicted from a school's public text alone using
two different methods.
The first is a prediction model: using the visited schools, a
statistical model learns which words and phrases of a document predict enacted scores, and it is validated only on schools that were not used to estimate it. Two things are done before estimation. First, identifying details are removed from the text. Second, each enacted score is regressed on three intake measures and the age of the inspection report, and the model predicts only the part of the score those variables do not explain. Without this step a model could predict the scores accurately just by recognising a school's intake from its report rather than by reading anything about its culture.

The second method uses a mark scheme approach. A large language model is provided with a one-page written mark scheme for each dimension --- warmth, strictness, and teaching --- and asked to score each document on a 1 to 5 scale. The first method assumes nothing about where
culture shows up in text; the second assumes it is there and scores it.

Each source is read on its own rather than pooled with the others. The three documents have two authors: the school writes its behaviour policy and its website, and an inspector writes the report. Pooling them would mix what a school says about itself with what an inspector says about it, and this paper asks about the two separately. There is a second reason. The language model applying a mark scheme reads one document at a time and gives it a band, so it cannot be given three documents at once. The prediction model is given one document at a time as well, so that both methods read the same input and their results can be compared. Pooling is reported later, for the prediction model on its own.

Each method gives every school a score. A set of scores is used for the national extension only if it agrees with what the observers recorded, and only if that agreement survives the checks described next. The two methods need different checks, because they can fail in different ways. A prediction model has many thousands of phrases to choose from, so it might reach some correlation by chance alone. A permutation test measures how much: the whole procedure is run again on scores shuffled between schools, which destroys any real link between a report and its score, and the real result has to beat what those shuffled runs produce. A mark scheme might succeed because of the particular language model applying it, or because it is following how long a document is rather than what it says, so its result has to survive a second scoring run, a different language model, and the removal of any effect of length. Across all the sources, methods and dimensions, many results are tested, so some would look convincing by chance alone. The significance levels are therefore corrected for the number of results tested, and a result counts only if it stays significant after that correction.

There are four main findings. The first is that inspection reports
carry a modest but real signal of enacted warmth. A prediction model reading only the report's wording predicts the enacted warmth of schools it was not estimated on. Its predictions correlate with what the observers recorded at $\PredWOfsted$. When the model predicts the raw enacted warmth score, the correlation is $\PredWOfstedRaw$. When it predicts only the part of the score that intake and report age do not explain, the correlation rises to $\PredWOfsted$. If the model were recognising intake from the report, removing intake from the target would lower the correlation. It does not, so the model's predictions rest on something other than intake. The phrases that lower a school's predicted warmth differ in kind from those that raise it. The reports of schools with low enacted warmth use stock assurances such as ``pastoral support is in place'' and ``pupils feel confident they have a trusted adult''. The reports of warmer schools instead describe pupils and adults
dealing with one another, spread across many small choices of phrase
rather than concentrated in quotable lines. The same is true of observed teaching. Predicting each school from a model estimated without it, the model's predictions correlate with the observers' teaching scores at $\PredTOfsted$. The mark scheme applied to the same reports does not predict warmth: its bands correlate with enacted warmth at $\RubricWOfstedDec$.

The second finding is that each method works for a different dimension: the prediction model for warmth, the mark scheme for strictness. The prediction model predicts enacted strictness poorly, at $\PredSOfsted$, where the mark scheme's bands correlate with it at $\RubricSOfstedDec$. On warmth the order is reversed. The reason is that the two dimensions appear in the text differently. Warmth is associated with a report's general vocabulary, which a model can learn from the visited schools without being told in advance what warmth looks like. Strictness is associated with specific factual claims, such as ``low-level disruption is rare''. The mark scheme states in advance what a strict school's report should say, so the language model applying it is told which sentences to look for. The prediction model is given only the reports and the observers' scores, and must find the association from the reports alone. For strictness it fails to.

The third finding is that the enacted and espoused measures diverge,
and that the divergence is stronger for warmth. The interview and the
visit are correlated at only $\GoldWarmthSplit$ for warmth, against $\GoldStrictnessSplit$ for strictness. Both measures contain random error. The error in one score is unrelated to the other score, so that part of each score cannot contribute to the correlation between them. The highest correlation the two scores could reach is therefore below one, and the raw figures understate how far the two agree. Correcting for the error in both raises warmth to $\DisSplitW$ and strictness to $\DisSplitS$. Two error-free measurements of the same quantity would correlate near one. Both corrected figures fall well short of that, so the head's account and the observers' scores are measuring related but distinct quantities.

The fourth finding is that warmth and strictness divide the same way again in what schools write about themselves. Models estimated on a school's behaviour policy and website predict what its head says about strictness, correlating at $\PredEspS$, and more weakly what the head says about staff morale, at $\PredEspClim$. On warmth they predict neither: $\PredEspW$ against the head's account, and no better against what the observers saw.

The paper contributes to various strands of the literature. The first is the school-climate literature, whose standard tool is the survey: pupils, teachers and parents are asked what their school feels like \citep{hoy2001,cohen2009,thapa2013}. Schools that come out better on these surveys have better wellbeing, attendance and attainment. But a survey records what people say about their school, not what the school does, and people tend to describe their own school kindly. \citet{vanhoutte2005} makes the distinction directly: surveys measure climate, the sum of what members perceive, while what shapes behaviour is culture, the shared norms of the school. It is culture in that sense that this paper measures. The cheaper alternatives to a survey do not measure it either. Exclusion rates \citep{machin2020} reflect who a school's pupils are as much as how the school treats them, and the inspection grade compresses everything a school does into a single verdict. What is missing is a measure of practice, taken by watching schools rather than asking them, that could in principle be extended to every school. That is what the paper builds.

The second is the school-effectiveness tradition that begins with \citet{coleman1966} and \citet{rutter1979}. Its question is why some schools do better than others with similar pupils. Its answer was not resources but what Rutter called ethos: the expectations a school holds and the practice it keeps up day to day. Later studies agreed on what a good ethos involves, and the list pairs an orderly environment with a supportive climate \citep{edmonds1979,sammons1995,teddlie2000}. That pair is close to the strictness and warmth measured here. But the tradition has studied ethos through case studies, one school at a time, so its central idea has never been measured in a way that lets schools be compared. The strictness and warmth scores do that for its two main ingredients.

Structured classroom observation, in particular the CLASS framework \citep{pianta2012}, measures emotional support and classroom organisation directly and predicts pupil gains at secondary level \citep{allen2013}; observation instruments from the same family predict gains in English primary schools \citep{sammons2008}. The
observation forms I use draw on those two domains, but they aggregate across
lessons and out-of-lesson settings so as to characterise a school rather
than a teacher. To my knowledge no study has combined such observation with
a leadership interview and public text in order to separate what a school
does from what it says.

The literature on high-expectation, high-structure
charter and academy models
\citep{abdulkadiroglu2011,angrist2010,dobbie2013,fryer2014} finds large
effects of a bundle of practices in which strict behaviour norms and warm,
individual relationships come together. It is the origin of the strictness
dimension and the reason to expect it to matter, but it has not measured the two dimensions apart.

Finally, the text-as-data literature in economics
\citep{gentzkow2019} supplies the tools for reading public text at
scale.  This paper uses both of that literature's modern tools, supervised text regression and mark schemes applied
by a large language model \citep{gilardi2023,eloundou2024}. It tests each against direct observation
and reports which dimension of culture each can read. \citet{ludwig2025} argue that a language model's measurements cannot be trusted on their own and must be checked against a gold standard. This paper carries out exactly that check for one concrete measurement problem, and goes one step further by running the statistical alternative against the same gold standard.

Against those literatures, the paper contributes three things. The first is a measure of enacted culture built from observation rather than perception; the perception data from headteachers and parents are kept as comparisons against it. The second is a model that predicts enacted culture in every English state secondary school. The third is a comparison of the two ways economists now read text, prediction models and mark schemes, with both judged against the same enacted scores for the same schools. The comparison shows that neither method is better outright: which one works depends on what it is being asked to measure.

The next section describes the data. The paper then constructs the gold standard, describes how the two text measures are built and judged, reports what the measures show, and concludes.

% ---------------------------------------------------------------
\section{Data}
\label{sec:p1_data}

The paper uses four kinds of data on English state secondary schools: a structured observation
from a visit and a
headteacher interview that were carried out for this
study, three collections of public text, a national survey of parents, and
the Department for Education's school-level administrative data. The sources divide by whose
account of the school they give.
The visit records what trained observers saw the school do, and the inspection report
records what an inspector saw, so these two are sources on \emph{enacted}
culture. The interview, the behaviour policy and the website record what the
school says about itself, the first to a researcher and the other two in
writing, so these three are sources on \emph{espoused} culture. The parent
survey belongs to neither group. Parents report their own family's experience
of the school, and whether that experience follows what observers see or
what the school says about itself is one of the questions the paper answers.

The study therefore has two samples. In \NTierOne{} schools there is a visit, an interview and the public text. In a further \NTierTwo{} schools there is an interview and the public text but no visit, so 303 schools carry an interview score in all. Headteachers gave verbal consent to the interview and the visit, and schools are not named. For the rest of the schools in the country there is the public text alone. \Cref{tab:tier_summary} compares the three groups on intake, size and inspection grade.

\begin{table}[htbp]
\centering
\small
\begin{tabular}{lcccccc}
\toprule
Group & $N$ & FSM\% & EAL\% & Median size & Outstanding\% & Good\% \\
\midrule
Visited and interviewed & 103 & 29.2 & 27.8 & 1{,}118 & 25 & 65 \\
Interviewed only & 200 & 27.2 & 17.9 & 1{,}052 & 15 & 75 \\
Rest of the national population & 3{,}028 & 30.1 & 18.2 & 1{,}050 & 14 & 70 \\
\midrule
All schools & 3{,}332 & 29.9 & 18.4 & 1{,}051 & 15 & 70 \\
\bottomrule
\end{tabular}
\caption{Distribution of schools across the three groups and selected observable characteristics (2023--24 school year). The groups are mutually exclusive: the first received both a headteacher interview and a researcher visit, the second an interview only, and the third is the rest of the national population. FSM\%: mean proportion eligible for free school meals.
EAL\%: mean proportion with English as an additional language. Size: median
number of pupils on roll. Outstanding\%/Good\%: percentage rated Outstanding or
Good at most recent inspection, among schools with a grade.}
\label{tab:tier_summary}
\smallskip
\begin{minipage}{0.92\linewidth}
\footnotesize\textit{Notes:} The groups sum to one fewer than the total: one interviewed school whose head left the statement battery blank belongs to the total but to no group.
\end{minipage}
\end{table}


\textbf{Visits and interviews.} Between November 2023 and April 2025 the
research team, working under my direction, interviewed the headteacher at
304 schools and visited \NTierOne{} of them. In eight cases the head was
unavailable and we interviewed the senior leader responsible for behaviour
and culture instead.\footnote{The counts that follow differ by a school or
two from place to place, for four reasons which I state once here. One of
the 304 heads left the statement battery blank, so 303 schools carry an espoused score and the interviewed-only group numbers \NTierTwo{}; that school belongs to neither group in \cref{tab:tier_summary}. Of the \NTierOne{}
visited schools, two have no scored inspection report and two have no
published parent survey, so correlations and prediction models that use
those sources run on 101 schools. The demographically adjusted correlations run on 100, because one school lacks a control variable.
Among the interviewed schools, between 299 and 302 have each text source. Finally, one visited school and three
interviewed schools could not be matched to the 2023/24 characteristics
panel used in \cref{app:2A:sample}, which therefore describes 102 and 300
schools.}

Schools were recruited by approaching headteachers directly by
email and telephone, and by referral from schools that had already taken
part. No quota was applied by region, school type or inspection grade.
Three-quarters of the interviews and half of the visits took place in the
2024/25 academic year. At the end of every interview the headteacher was
asked whether the school would be willing to host a visit, which means that
the visited schools are a self-selected subset of the interviewed ones.

A visit lasted a school day and was made by a team of two, and occasionally
three, researchers. The school was asked in advance to arrange four to six
lessons spread across year groups and subjects, and it chose the particular
lessons itself. The team observed those lessons, typically five of them,
spread across Years 7 to 11. The researchers knew which school they were
visiting but not its inspection grade, and the school knew it was being
observed but not which items were being rated. Each school also allowed us
to observe pupils moving between lessons, and pupils and staff interacting
at break and lunch, which is what the outside-of-lessons form rates. Each observation covered
either a whole lesson or half of one, and every researcher present rated it independently on the
lesson observation form. The team also watched arrival to lessons, movement between lessons,
break and lunch, and
at the end of the day each researcher completed the outside-of-lessons form once,
summarising the whole day.

The lesson form has twenty items; the mapping from the printed form to the scored items
is in \cref{app:2A:guide}. Ten of the
twenty relate to warmth and strictness, for example the
frequency with which the teacher uses pupils' names, and the level of
misbehaviour tolerated before a first sanction. The other ten relate to
teaching practice, such as questioning, feedback, explanation and structure.
The outside-of-lessons form has thirteen items: nine on the school outside
the classroom (sanctions and rewards, corridors, arrival, the dining hall,
and interactions and relationships at break), two on the condition of wall
displays, and two in which the observer compares the teaching seen across
the day. Every item is rated from 1 to 5 and the item ratings are averaged across
researchers and across lessons to give one value per school.

The observers saw the lessons the school chose to show them. What they scored may therefore be closer to the school's best lessons than to its typical ones. But the likely direction of the problem is known. A school will tend to show a visitor its better lessons. If every school does this, every school's enacted score is a little too high, and the scores are more alike than the schools really are. Scores squeezed together make differences between schools harder to find, not easier, so this bias works against the paper's findings rather than for them. One version of the problem remains open: some schools may be better at putting on a show for a visitor than others. That difference could push the scores in either direction, and nothing in the data rules it out.

The interview lasted about an hour and followed a fixed guide. It had three
parts. The first part consisted of twelve open questions on the school's approach to behaviour,
relationships, wellbeing and teaching.
This was followed by ten statements on which the
headteacher rated their own school on a five-point agreement scale (for
example ``The sanction systems are applied consistently by all staff'' and
``Pupils feel like their teachers care''). The final part was a checklist of ten named behaviour systems (centralised detentions, an
internal exclusion room, a mobile phone ban and so on) recorded as present
or absent. Two researchers attended every interview, one asking the
questions and one observing, and both scored the open questions
independently. The observation forms and the interview guide are reproduced
in full in \cref{app:2A:visit_protocol,app:2A:interview_guide}, and the
agreement between researchers on the observation and interview items is reported together
with the measures below.

\textbf{Public text.} Three collections of text cover the schools in the national population. A scored inspection report exists for
3,272 of them (98.2 per cent); the rest had opened too recently to have been
inspected. A behaviour policy document exists for 3,294 (98.9 per cent),
one document per school, taken from the school's website. About one document
in ten is a trust-wide policy adopted word for word by the school. Website
text exists for 3,284 schools (98.6 per cent). This text is taken from the home,
values and welcome pages. The counts differ by a few dozen schools because a
handful of schools have no report, no policy or no website;
\cref{app:2A:sample} gives the reconciliation.

\textbf{Parent survey.} Ofsted's Parent View survey is open to the parents
of every school and is published once enough parents have responded. I use
the most recent release for each school and the two items closest to
warmth: the share of responding parents who agree that their child is happy
at the school, and the share who agree that their child feels safe there.
The two shares are averaged into one parent-reported warmth measure, which
is available for 3,255 schools in the panel and for 101 of the visited
schools. Parents choose whether to respond, and the typical release rests on about 120 responses at a response rate of about 12
per cent, so the measure is the view of the parents who chose to give one. Responses also accumulate over several years,
so for some schools the parent data are older than the culture data. The survey has no item on strictness.

\textbf{Administrative data.} All schools are matched on their unique
reference number to a school-year panel that I built from the Department for
Education's published tables for 2015/16 to 2024/25. The panel carries pupil
composition, prior attainment, Progress~8 (a value-added measure) and its components,
workforce and
finance. This paper uses the panel only to describe the samples; the
following chapter uses it for outcomes and controls. The panel also carries
each school's most recent overall inspection grade. Inspections from
September 2024 onwards no longer award an overall grade, so the grade used
throughout is the last one each school received before that date, on the
scale 1 (Outstanding) to 4 (Inadequate). Schools inspected only since then,
or never inspected, have no grade and are missing on that variable.

\textbf{Representativeness.} \Cref{tab:2A:representativeness} compares the
visited and interviewed schools with the national population. Both samples
resemble the population on school type, phase, admissions, size and urban
location, and they differ from it on deprivation by at most two percentage
points. Both over-represent London (45 per cent of visited schools against
15 per cent nationally), which shows up as a higher share of pupils with
English as an additional language. Both also over-represent stronger schools:
25 per cent of visited schools held an Outstanding inspection grade as at
August 2024 against 15 per cent nationally, and their mean Progress~8 in
2018/19 was $+0.17$ against $+0.04$. The mismatch on those dimensions is what one would expect from recruiting through headteachers,
since the schools most likely to agree to take part are those that are
confident about what a visitor will find, and it has two consequences. The first is that
estimates from the visited sample apply most directly to a population that is somewhat stronger than the national average.
The second is a restricted range: the sample under-represents the weakest schools, so both culture and outcomes
vary less than nationally.

% ---------------------------------------------------------------
\section{Measurement}
\label{sec:p1_measurement}

This section builds the measures starting with the gold-standard scores from the visits and
the interviews, and reports how reliable they are. It then builds the prediction model and describes
the mark scheme and LLM approach that were both used to
produce scores for the schools that were not visited. It ends with how both approaches are judged.

\subsection{The gold standard}

The scores are built up in two steps. Every observation item
and every interview statement is rated from 1 to 5. First, items are averaged into
\emph{sub-scores}, each of which covers one dimension in one setting: warmth
in lessons, warmth outside lessons, strictness in lessons, strictness outside
lessons, and so on. Second, the sub-scores for a dimension are averaged into
that dimension's \emph{score}, so that each school ends up with one warmth score and one strictness score on the same 1--5 scale as the items.

Two sets of scores are constructed and kept apart. The \emph{enacted} scores use the observations
alone. Enacted
warmth is the average of in-lesson warmth (the use of names, praise, the
quality of interaction, pupil participation and motivation) and
out-of-lesson warmth (relationships and interactions during transitions and
breaks, and how consistently rewards are used). Enacted strictness is the
average of in-lesson strictness (misbehaviour tolerated, disruption,
response to sanction, respect and concentration) and out-of-lesson
strictness (how consistently sanctions are used, corridors, arrival, the
dining hall and recreational space). Enacted teaching is the average of the
ten lesson items on teaching practice.

The \emph{espoused} scores use the
headteacher's statement ratings alone. Espoused warmth is the average of the
three warmth statements (rewards applied consistently; pupils feel safe;
pupils feel their teachers care) and espoused strictness is the average of
the four strictness statements (sanctions applied consistently; behaviour in
lessons; behaviour outside lessons; teachers feel respected). The two
remaining scored statements, on staff turnover and morale, form a
staff-climate score. The item definitions and the mapping of items to
sub-scores are given in \cref{app:2A:guide}.

To make the construction concrete, consider the in-lesson strictness
sub-score of a single school. In each of five lessons, two researchers rated
the level of misbehaviour tolerated before a sanction, the frequency of
disruption, the pupils' response to correction, how respectfully they treated
the teacher, and how well they concentrated. Misbehaviour and disruption are
reverse-scored so that a higher value means a stricter lesson.

One item needs a rule: if no pupil was sanctioned during a lesson there is no response
to correction to rate, and the researcher records a zero. That zero is not a
low score, and dropping it would throw away what the lesson shows. It is
therefore replaced by the reverse of the disruption rating, which means that
a lesson with no disruption and no sanction is scored as if the pupils had
responded fully, and a lesson with constant disruption and no sanction as if
they had not responded at all.

The five ratings a researcher makes in a
lesson are averaged to give that researcher's value for the lesson, the
lesson values are averaged
across lessons, and the researchers' values are averaged across researchers,
giving one in-lesson strictness figure on the 1--5 scale. That figure is averaged with the out-of-lesson strictness figure to give the school's enacted strictness.

Espoused strictness is
simpler: the headteacher's ratings on the four strictness statements are averaged. A similar process produced espoused warmth.
Every school that received a visit has an enacted
score, every school that gave an interview has an espoused score, and
\NTierOne{} schools have both. \Cref{tab:score_dist} reports the distribution
of every sub-score and score across the visited schools.

\begin{table}[htbp]
\centering
\caption{Distribution of sub-scores and of the enacted and espoused scores for
Tier~1 full-data schools ($n = 103$).
Sub-scores are on a $[0,5]$ scale; the enacted and espoused scores on a $[0,10]$ scale.}
\label{tab:score_dist}
\footnotesize\setlength{\tabcolsep}{4pt}
\begin{tabular}{llcccccc}
\toprule
Score & Description & $N$ & Mean & SD & Min & Median & Max \\
\midrule
$W1$ & In-lesson staff--pupil warmth & 103 & 3.24 & 0.45 & 1.92 & 3.22 & 4.47 \\
$W2$ & Warmth in transitions and breaks & 103 & 3.32 & 0.68 & 1.83 & 3.33 & 4.67 \\
$W3$ & HT espoused warmth (statements) & 103 & 4.20 & 0.47 & 2.33 & 4.33 & 5.00 \\
$S1$ & In-lesson behaviour management & 103 & 3.69 & 0.49 & 2.12 & 3.70 & 4.75 \\
$S2$ & Out-of-lesson behaviour management & 103 & 3.33 & 0.50 & 2.10 & 3.30 & 4.60 \\
$S3$ & Behaviour management systems (count) & 103 & 3.11 & 0.73 & 0.62 & 3.12 & 4.38 \\
$S4$ & HT espoused strictness (statements) & 103 & 3.95 & 0.58 & 1.75 & 4.00 & 5.00 \\
$T1$ & In-lesson teaching quality & 103 & 3.48 & 0.42 & 2.49 & 3.51 & 4.51 \\
$SC$ & HT espoused staff climate (statements) & 103 & 3.91 & 0.78 & 1.50 & 4.00 & 5.00 \\
\midrule
Warmth & \emph{enacted} & 103 & 6.56 & 0.99 & 4.33 & 6.60 & 8.81 \\
Strictness & \emph{enacted} & 103 & 7.02 & 0.87 & 4.85 & 7.04 & 8.96 \\
Teaching & \emph{enacted} & 103 & 6.97 & 0.83 & 4.98 & 7.02 & 9.02 \\
Warmth & \emph{espoused} & 103 & 8.39 & 0.95 & 4.67 & 8.67 & 10.00 \\
Strictness & \emph{espoused} & 103 & 7.90 & 1.16 & 3.50 & 8.00 & 10.00 \\
Staff climate & \emph{espoused} & 103 & 7.81 & 1.55 & 3.00 & 8.00 & 10.00 \\
\bottomrule
\end{tabular}
\begin{minipage}{\textwidth}\vspace{0.5em}\footnotesize
\textit{Notes.} The enacted scores are built from the visit sub-scores
($W1$, $W2$, $S1$, $S2$, $T1$) and the espoused scores from the interview
statement sub-scores ($W3$, $S4$; $SC$ is the staff-climate score), as set
out in \cref{eq:enacted,eq:espoused}. The systems count $S3$ enters no
score. No school contributes a sub-score to both sets.
\end{minipage}
\end{table}


There are two larger exclusions from the interview. The first is the ratings of how fully the headteacher answered each open question. They correlate with the enacted scores at only $\ArticEnW$ for warmth, $\ArticEnS$ for strictness and $\ArticEnT$ for teaching, and no more strongly with the head's own ratings ($\ArticEspW$ and $\ArticEspS$). They therefore appear to record how fully the head talks rather than what the school is like, and they are excluded. The second is the count of the named behaviour systems (S3 in the tables), which counts eight of the ten and excludes the two reward systems. It correlates with enacted strictness at $\SystemsEnS$. In other words, the systems a school names are not a guide to how strictly it runs.

Agreement between researchers is moderate at the level of a single item, and
the school-level scores built from many items, lessons and raters are
reliable enough to use, as I show below. Of the 517 lessons observed,
\LessonsTwoRaters{} were rated independently by two researchers and
\LessonsThreeRaters{} by three; in the remaining \LessonsOneRater{} only one
researcher was available for that lesson, and it carries one rating. The
outside-of-lessons form was completed by two or more researchers on every
one of the \NTierOne{} visits. A researcher left an item blank when they
did not see what it asks about, so for any one item the number of visits
with two ratings to compare is smaller than \NTierOne{}.

Where two raters saw the
same lesson, the pairwise Pearson correlations across the classroom items
range from $0.53$ to $0.79$; on the outside-of-lessons form, completed once
per day, they range from $0.47$ to $0.81$; and the two raters are within one point of each other close to 90 per cent of the time. \Cref{tab:irr_classroom,tab:irr_outside} report the agreement item by item. Agreement is highest on the
items whose scale points are described most concretely on the form, namely the consistency of the reward
system ($0.81$, a day item) and the frequency of disruption ($0.79$), and
lowest among the classroom items on the teaching-practice ratings that ask
the observer for the most interpretation, such as the use of resources and
the choice of methods ($0.53$ and $0.55$).

However, the paper's claims do not rest on the agreement of two raters on
one lesson. They rest on the reliability of a school's score once it has
been averaged over its lessons and raters. \Cref{tab:score_reliability} gives this. Take
every rating that a
researcher gave a lesson for in-lesson strictness: one number, for one
lesson, from one researcher. Those ratings differ from one another for three
reasons. The first is that the lesson was in a stricter or a less strict
school, which is the difference the measure is meant to capture. The second
is that, within the same school, that lesson was calmer or rowdier than the
school's other lessons, because it had a different teacher, a different
class or a different period of the day. The third is that, looking at the
same lesson, one researcher scored it a little higher than the other.

The table divides the spread of the ratings into those three shares. For
in-lesson strictness the shares are \ShareSchoolS{} per cent school,
\ShareLessonS{} per cent lesson and \ShareRaterS{} per cent rater. This means
that well over half of the variation in single-lesson strictness ratings is
lessons within the same school differing from one another, only about a
fifth is schools differing from each other, and under a fifth is the two
researchers disagreeing. For in-lesson warmth the shares are
\ShareSchoolW{}, \ShareLessonW{} and \ShareRaterW{} per cent. Lessons within
a school vary less in warmth than in strictness, but the researchers
disagree more about warmth, which matches the item-level pattern described
above.

\begin{table}[htbp]
\centering
\small
\begin{tabular}{lccccc}
\toprule
 & \multicolumn{3}{c}{\textit{Share of variance (\%)}} & & \\
\cmidrule(lr){2-4}
Sub-score & School & Lesson & Rater & Design & Reliability of school mean \\
\midrule
In-lesson warmth ($W1$) & 21 & 38 & 42 & 4.7 lessons $\times$ 1.6 raters & 0.60 \\
In-lesson strictness ($S1$) & 21 & 61 & 18 & 4.7 lessons $\times$ 1.6 raters & 0.58 \\
Teaching practice ($T1$) & 27 & 40 & 33 & 4.7 lessons $\times$ 1.6 raters & 0.67 \\
\addlinespace[3pt]
Out-of-lesson warmth ($W2$) & 79 & --- & 21 & 1.6 raters & 0.85 \\
Out-of-lesson strictness ($S2$) & 71 & --- & 29 & 1.6 raters & 0.79 \\
\bottomrule
\end{tabular}
\caption{Variance components and score reliability}
\label{tab:score_reliability}
\begin{minipage}{0.92\linewidth}
\vspace{4pt}\scriptsize
\textit{Notes}: Lesson sub-scores are decomposed into between-school, between-lesson-within-school and between-rater-within-lesson components; the outside-of-lessons sub-scores, rated once per rater per school, into between-school and between-rater. Reliability is the share of variance in a school's mean score, over its own lessons and raters, that is between schools. Method-of-moments components from nested one-way analyses of
variance on 517 lesson observations in 103 schools
(lesson sub-scores) and 103 schools (outside-of-lessons sub-scores);
design sizes are harmonic means. Reliability is $\sigma^2_{\mathrm{school}} /
(\sigma^2_{\mathrm{school}} + \sigma^2_{\mathrm{lesson}}/m_l + \sigma^2_{\mathrm{rater}}/m_l m_r)$.
The in-lesson and out-of-lesson halves of the enacted score, rated in different
settings, are correlated at $r = 0.52$ for warmth and $0.53$ for
strictness across schools.
\end{minipage}
\end{table}


One implication of dividing the variation into those three shares is that averaging matters. A single lesson is a poor guide to its school, but when
five lessons and the ratings of the researchers who saw them are averaged,
the lesson-to-lesson and rater-to-rater variation is greatly reduced. The reliability of the
in-lesson warmth score,
averaged over a school's own lessons and raters, is \RelWLessons{}, and that of the in-lesson
strictness score is \RelSLessons{}. The outside-of-lessons sub-scores, which summarise a whole day,
are more
reliable (\RelWOutside{} and \RelSOutside{}).

Each enacted score combines two scores: a lesson score, from the observed
lessons, and an outside-of-lessons score, from the rest of the school day.
The two are correlated at \HalvesW{} for warmth and \HalvesS{} for strictness across
the \NHalves{} schools. Treating them as two readings of the same score, the reliability of the combined score is that correlation stepped up for the doubling of length. This puts the reliability of each combined enacted score at
about \RelEnactedW{} for warmth and \RelEnactedS{} for
strictness. The reliability is the share of the score's variance that is not error. A score correlates with its own error-free value at the square root of that share, and no other measure can correlate with the score more highly than that. For enacted warmth this limit is $\sqrt{0.69} \approx 0.83$.

These figures rest on an assumption about why the lesson score and the
outside-of-lessons score differ. They do not agree perfectly, and the
figures above treat all of that disagreement as
measurement error. That is the right treatment if lessons and the rest of the day are two
readings of one and the same thing. But a school might genuinely be warmer in its lessons
than in its corridors. If so, part of the disagreement is real rather than error, and the
reliabilities are \RelEnactedWUpper{} and \RelEnactedSUpper{} instead. The corrections below
use the lower figures. Dividing by a lower reliability gives a larger corrected correlation,
so every corrected correlation in this chapter is the most generous version of that number:
the true agreement could be lower than the corrected figure, but not higher.

One further limitation: each school was visited once. If a school is different on a different day,
that difference cannot be separated from the school itself. These are reliabilities of
the school as seen on its visit day.

The lessons that carry only one rating lower the reliabilities very little. Had every lesson been rated twice, the in-lesson
reliabilities would have been \RelWOneTwoRaters{} and \RelSOneTwoRaters{}
rather than \RelWLessons{} and \RelSLessons{}. This is because most of the variation within a school is between lessons, not between raters, as the shares above show. A second rater on every lesson would reduce only the smaller of the two. They do, however,
carry one caution. Single-rated lessons were rated lower on average than double-rated ones. The difference, single minus double, is $\SingleRatedGapW{}$ on warmth
and $\SingleRatedGapS{}$ on strictness on the 1--5 scale. The \SchoolsWithSingleRated{} schools with such lessons have
slightly lower scores as a result: the share of a school's single-rated lessons correlates with
its in-lesson warmth at $r = \CorrSingleShareW{}$. The following chapter therefore checks that its
estimates do not depend on those schools.

The division of the variation into a school share, a lesson share and a rater share also says
something about the schools themselves, not just the measurement. Behaviour varies more from
classroom to classroom within a school than it does between schools. A school's enacted score
is therefore an average over
quite different lessons, and a visit of five lessons measures that average with error. Random error in a score adds to its variance without adding to its covariance with any other measure. Every correlation with the enacted scores is therefore lower than it would be without error, so the figures reported below are, if anything, too small.

Two corrections give a sense of how much the error lowers the correlations. The first correction divides the raw enacted--espoused correlation by the square root of the product of the two scores' reliabilities. The result is the correlation the two scores would have if neither contained error. It raises warmth to about \DisSplitW{} and strictness to about \DisSplitS{}. Two error-free measurements of the same quantity would correlate near one, so even after the correction the head and the observers are far from describing the same school. The second corrects the mark scheme's strictness correlation with the enacted score, $\RubricSOfstedRaw$, for the enacted score's error alone, and it rises to about $0.36$. In both cases the corrected number leads to the same verdict as the raw one, so the uncorrected figures are the ones reported throughout.

The two enacted scores are correlated at $r = \CorrWarmthStrictnessEnacted$ across the \NTierOne{} visited schools, and the two espoused scores at $\CorrWarmthStrictnessEspoused$. Taken at face value, the enacted figure says that warm schools are usually strict ones. Part of that figure, though, comes from how the scores were made rather than from the schools. The same observers scored both dimensions at the same school on the same day, and an overall impression of a school formed on the day can move both ratings in the same direction.

The lessons that both researchers rated independently show where the correlation comes from. Suppose it arose within each observer, with the same person's judgement of warmth affecting their judgement of strictness. Then an observer's warmth rating would predict their own strictness rating much better than their colleague's. The two correlations are nearly the same: \CorrWSSameObs{} and \CorrWSCrossObs{}. So the correlation reflects what both observers saw rather than the judgement of either one. The cross-observer figure still contains ordinary measurement error. Dividing it by the square root of the product of the two enacted reliabilities gives about \CorrWSTrue{}. That, rather than $\CorrWarmthStrictnessEnacted$, is the better estimate of how strongly warmth and strictness go together in schools. The two are related but clearly distinct.

The gold standard is two measures rather than one because the interview and the visit disagree about
the same school. If the two were noisy readings of one quantity, the right thing to do would be
to average them,
because their errors would partly cancel. However, they agree far too little for that: the
two warmth measures correlate at about 0.25, so each explains about one sixteenth of the other's variance, and the two strictness measures at about 0.41, so about one sixth. Averaging two measures of different things does not cancel noise. The average would be a mixture, partly what the observers saw and partly what the head said, and choosing different weights would give a different mixture.

I therefore keep
the two measures separate throughout, and their divergence is itself a result. For strictness the corrected correlation reaches \DisSplitS{}. The head's account and the observers' scores therefore share a substantial part of their variance, but well under all of it. For warmth the corrected correlation is \DisSplitW{}, so the two share little. Two schools make this concrete, one where the two measures agree and one where they diverge. At the first, the headteacher rated the school 4.0 on warmth and 4.0 on strictness, and the observers, across their lessons and the day, recorded 3.9 and 4.3; the in-lesson and out-of-lesson sub-scores agree with one another to within half a point. At the second, the headteacher rated warmth 4.7 and strictness 3.9, near the top of the scale, while the observers recorded 2.3 and 3.0. In-lesson warmth at that school
averaged 2.2 out of 5 and out-of-lesson warmth 2.4, while the head's ratings
on the three warmth statements averaged 4.7 out of 5. The two interviews read alike: both heads placed their schools high on the same warmth statements. Only the visit separates them. Across the \NTierOne{} schools, four in five have a warmth gap larger than half a point and half have a gap larger than one point.

One espoused measure can be checked against an outside record. Heads rated how low their staff turnover was and how high
morale was, and the staff-climate score built from those two statements is
correlated with the school's actual teacher turnover, as recorded by the
Department for Education, at $r = -0.52$. In a regression of teacher turnover on the staff-climate score, the other espoused scores, disadvantage and school size, the staff-climate coefficient corresponds to a partial correlation of $-0.48$. On the one
question where a head's answer can be checked against the record, heads
report accurately. This check concerns staff turnover, though, and says nothing about
the warmth statements.

\subsection{Predicting the scores from text}
\label{sec:p1_prediction}

The first way of reading the documents uses no theory of culture at all. It asks only whether
the words a document uses tend to go with the culture observers found in the school.

The method works as follows using warmth as an example. Each document is turned into a
list of counts:
how often each word, and each two-word phrase, appears in it. Phrases that
appear in almost every document count for little; phrases that distinguish
one document from another count for more. A statistical model is then estimated on the visited schools. It learns which phrases tend to appear in
the reports of schools observed to be warm, and which in the reports of
schools observed to be cold, and it attaches a small positive or negative
weight to each. To score a new school, the model adds up the weights of the
phrases in its document. No single phrase decides anything. A school's
score is the sum of thousands of small weights, each attached to one phrase. In symbols, let $x_s$ be school $s$'s text vector, $y_s$ its enacted
score and $C_s$ its four control variables. The target is the residual
from the demographic adjustment,
%
\begin{equation}
  \tilde{y}_s = y_s - \hat{c} - C_s'\hat{\gamma},
  \label{eq:p1_deconf}
\end{equation}
%
with $\hat{c}$ and $\hat{\gamma}$ estimated by least squares on all
schools. For each school $s$, the model is estimated on every other school,
%
\begin{equation}
  (\hat{a}_{-s}, \hat{\beta}_{-s})
  = \arg\min_{a,\,\beta}
    \sum_{t \neq s} \bigl(\tilde{y}_t - a - x_t'\beta\bigr)^2
    + \lambda_{-s}\,\lVert\beta\rVert^2,
  \label{eq:p1_ridge}
\end{equation}
%
with the penalty $\lambda_{-s}$ chosen by leave-one-out error within the
$t \neq s$ schools alone. The held-out prediction for school $s$ is
$\hat{y}_s = \hat{a}_{-s} + x_s'\hat{\beta}_{-s}$, and every correlation
reported for a prediction model is the correlation between $\hat{y}_s$
and $\tilde{y}_s$ across the held-out schools.
\Cref{sec:p1_app_prediction} gives the text representation and the remaining estimation detail.

Three safeguards stop the model from predicting the enacted scores accurately for the wrong reasons.

First, the model is never tested on a school it learned from. Every school's reported
prediction comes from a model estimated with that school left out, so each number says how accurately the model predicts the score of a school that was not used to estimate it.
As a further check, I ran a permutation test, re-running the whole procedure two hundred times on deliberately scrambled scores. Scrambling means shuffling the warmth scores across the schools, so that each report is paired with a score that belongs to some other school. The reports are real and the scores are real, but the pairing between them is destroyed, so there is nothing true left for the model to find. Whatever correlation it still reaches is what luck alone produces. A real result has to beat that, and the reported significance levels say how often the scrambled runs did.

The second safeguard is to ensure it cannot recognise the school. Before anything is counted, school
and trust names, place names, dates, page numbers and standard boilerplate
are removed from every document. Without this step, the model learns to spot which school or which era of report it is reading,
rather than what the report says.

Finally, the model is prevented from reading demographics. A school's intake
statistics predict part of its culture score, and a document often reveals
the intake. So each enacted score is first split into the part that
intake statistics explain and the part they do not, and the model is asked
to predict only the second part. A model that merely recognised a
disadvantaged school from its report would predict none of the adjusted score, so its correlation would be near zero. Results are
shown with and without this adjustment.

I also tried more elaborate ways of representing the text:
longer phrases, topic groupings, and meaning-vectors from a neural network that scores passages as
similar when they say similar
things in different words. \cref{sec:p1_app_prediction} records the findings. None of them improved on counting, so that is the method reported throughout. The more complex methods
failing to beat
counting does not cast doubt on counting. The complex methods can read things counting cannot.
A meaning-vector, for example, treats two sentences as alike when they say the same thing in
different words. So if the reports held a signal that counting misses, the complex methods
would have found it. The fact that they did not suggests that the signal in these documents is in
the words themselves.

\subsection{Scoring the text with mark schemes}
\label{sec:p1_rubrics}

The second way of reading the documents states what to look for. For each
source and dimension I wrote a mark scheme of one page: a sentence
saying what is being measured, five numbered bands, and two or three
tie-break rules. To give the flavour, the inspection-report warmth scheme
sets band~3, the default, as ``a safe and generally positive school,
described through labels, provision or outcomes, without describing how
staff and pupils actually relate'', and band~4 as ``the report describes
the relationships themselves''. All seven schemes are printed in full in \cref{sec:p1_app_rubrics}.

In the same notation, a mark-scheme band is $B_s = f(x_s)$. Here $f$ is the fixed procedure of giving the one-page scheme and one document to a named version of the language model and taking the band it returns. The scheme and the model were fixed before any document was scored. Nothing in $f$ is estimated from the schools, so there is no training step and no held-out distinction; the adjusted figure reported for a scheme is the correlation between $B_s$ and $\tilde{y}_s$.

A language model reads each document against the scheme and returns a band from 1 to 5, together
with one quoted sentence as
evidence (the models are named in \cref{sec:p1_app_rubrics}). The
same fixed model reads every document.

Three checks are carried out to assess the robustness of the LLM scoring, and their outcomes
are reported with the results. The first asks whether the language model gives the same answer
twice: thirty documents per source were scored a second time, and the agreement between the two
runs is reported. The second asks whether the bands depend on which large language model produced them. To assess this the whole exercise was repeated with a second, larger model. This
check also settles a tie. When the prediction model and the mark scheme correlate with the enacted scores about equally, the mark scheme is used only if its correlation is about the same under both language models. If the correlation changes substantially with the larger model, the prediction model is used instead.

The third check asks whether the mark schemes are reading length rather than culture. A longer document
gives the model more evidence to cite, so a scheme could award higher bands to longer documents
rather than to warmer or stricter schools. The correlation between band and document length is
therefore computed for every source and dimension, and each correlation is then recomputed after removing from both the band and the enacted score the part that document length predicts.

\subsection{Judging both methods}
\label{sec:p1_judging}

Both methods are judged in the same way: a method is good if its scores agree with what
trained observers recorded in the school, and it earns no credit for anything else.
Comparing the two methods source by source and dimension by dimension then answers two questions
at once. The first is how much the public text says about a school's culture at all. The second
is what it takes to read it: whether the words must be interpreted against a stated theory, as
the mark schemes do, or whether counting them is enough, as the prediction models do.

One further rule applies throughout. Warmth, strictness and teaching are different things, and
a school's three scores need not agree. So a signal that genuinely reads one of them should not
help predict all three at once. When a feature of a document does help with all three, the likely
reason is not that the school is better at everything. It is that the feature describes the
document itself, so two schools with similar documents receive similar scores however much the
schools differ. A signal of that kind is discarded. Document age, format and length were
discarded for exactly this reason, and so were several changes that appeared to improve
the models (\cref{sec:p1_app_prediction}).

\section{Results}
\label{sec:p1_results}

The results come in three parts: what direct measurement shows on its own, what the public
text contains and which method extracts it, and two checks from outside the project's own data.

\subsection{What direct measurement shows}

Direct measurement shows two things. The first is that warmth and strictness are
separate dimensions of a school. A school can be warm without being orderly, and
orderly without being warm, and the visited sample contains both kinds: \cref{fig:p1_quadrants}
plots the schools on the two dimensions and every quadrant is occupied. The two dimensions
are related but distinct: after the correction for shared observers and for measurement error, they correlate at \CorrWSTrue{}.

\begin{figure}[htbp]
  \centering
  \caption{Observed warmth and strictness in the visited schools}
  \label{fig:p1_quadrants}
  \includegraphics[width=0.8\textwidth]{fig_quadrants_ch2.pdf}
\begin{minipage}{\linewidth}\smallskip
\footnotesize\textit{Notes:} Each dot is a school. Dashed lines mark the sample medians, which divide the schools into the four labelled quadrants. All four are occupied: warm and strict; warm, less strict; strict, less warm; and neither.
\end{minipage}
\end{figure}

The second is that what heads say and what observers see are different measures. The interview and
visit scores are correlated at only $\GoldWarmthSplit$ for warmth ($p = \GoldWarmthSplitP$) and
$\GoldStrictnessSplit$ for strictness ($p \GoldStrictnessSplitP$), on the \NTierOne{} schools
with both.

Both scores are measured with error, so the highest correlation they could reach is below one. Correcting for the error in both raises warmth to $\DisSplitW$ and strictness to $\DisSplitS$. Neither reaches the value two error-free measurements of the same quantity would give, and the shortfall is much larger for warmth: heads place their school's strictness with some accuracy and its warmth with very little. The disagreement on warmth is therefore real, not a product of measurement error.

\subsection{What the public text contains}

This section shows three things. First, scores derived from public text do correlate with the enacted scores, but only when the text is an inspection report: every text measure that passes the tests comes from an inspection report, and the documents schools write about themselves carry no signal. Second, within the reports the two dimensions are read by different methods: the prediction model predicts warmth, the mark scheme scores strictness, and both work for teaching. Third, the measures the chapter goes on to use pass the three checks set out in the measurement section.

\Cref{fig:p1_validation} plots the warmth scores from the prediction model and the strictness
scores from the mark scheme against the enacted scores: the warmth model's held-out predictions
on the left, the strictness scheme's bands on the right. In both panels the two rise together:
schools the model reads as warmer are schools the observers scored as warmer, and reports the
scheme places in higher bands come from schools the observers scored as stricter. The relationship
is loose, with exceptions at every level, but it is not the work of a few outlying schools.

\begin{figure}[htbp]
  \centering
  \caption{The warmth model and the strictness scheme against the enacted scores}
  \label{fig:p1_validation}
  \includegraphics[width=\textwidth]{fig_ch2_validation.pdf}
\begin{minipage}{\linewidth}\smallskip
\footnotesize\textit{Notes:} Panel A shows the warmth prediction model's held-out predictions against enacted warmth, both after the demographic adjustment, that is, the part of each score that the four control variables do not explain; each school is predicted by a model that never saw it. Panel B shows the strictness mark scheme's bands against enacted strictness without that adjustment, so the correlation shown is the raw $\RubricSOfstedRaw$; the adjusted figure, $\RubricSOfstedDec$, is the one in \cref{tab:p1_grid}.
\end{minipage}
\end{figure}

\Cref{tab:p1_grid} reports fourteen text measures: seven pairings of a source with a dimension, each scored twice, once by the prediction model and once by the mark scheme. Each measure is summarised by one number: the correlation between the score it gives a school and the score the observers gave it. Each measure then faces two tests: it must beat luck, and it must stay significant once the number of tests is allowed for. Every measure that passes both tests comes from an inspection report, and none of the measures taken from the documents schools write about themselves passes either. Within the reports, warmth is read by the prediction model and strictness by the mark scheme. The rest of this section works through the tests in \cref{tab:p1_grid} and then the three checks, one at a time.

% generated by make_grid_numbers.py — do not edit
\begin{table}[htbp]
\centering
\caption{What each public text carries about enacted culture: two readings,
one criterion. Cells are correlations with the visit scores on the visited
schools ($n \approx 100$), with intake and report age partialled out of the
target. ``Agree'' is exact band agreement between the two language models
reading the same marking scheme; ``Retest'' is exact agreement on thirty
documents re-scored by the primary model.}
\label{tab:p1_grid}
\begin{tabular}{llcccc}
\toprule
Source & Dimension & Prediction model & Marking scheme & Agree & Retest \\
\midrule
Inspection report & Warmth & +0.33 & -0.07 & 68\% & 97\% \\
Inspection report & Strictness & +0.19 & +0.25 & 71\% & 100\% \\
Inspection report & Teaching & +0.28 & +0.27 & 72\% & 97\% \\
Behaviour policy & Warmth & +0.15 & +0.04 & 62\% & 83\% \\
Behaviour policy & Strictness & -0.15 & -0.03 & 72\% & 100\% \\
Website & Warmth & +0.14 & -0.13 & 31\% & 97\% \\
Website & Strictness & +0.17 & +0.13 & 39\% & 97\% \\
\bottomrule
\end{tabular}
\end{table}


The scrambled scores say how easily each of the model's results could be luck. Its three inspection-report measures all beat luck (permutation test, warmth \PermPOfstedW{}, teaching \PermPOfstedT{}, strictness \PermPOfstedS{}). Its measures from the behaviour policies and the websites all fail to.

The prediction model can be given several documents at once and the mark scheme cannot, so these pooled results are for the model alone. Each figure is again the correlation between the model's predictions and the observers' scores. Pooling the behaviour policy and the website, the two documents the school writes, gives a correlation with enacted warmth of $\PredBPWebW$; for strictness and teaching the pooled correlations are near zero. So pooling adds no signal that the separate readings missed. It can also lose one: the website on its own correlates with enacted strictness at $\PredSWeb$, and adding the policy leaves $\PredBPWebS$. Pooling all three sources never improves on the best single source; on observed teaching the correlation falls from $\PredTOfsted$ for the inspection report alone to $\PredTCombined$ for the three together.

Fourteen measures also give luck fourteen chances, so the measures in \cref{tab:p1_grid} are corrected for the number of tests as well. Four pass this test: the model's warmth and teaching, and the scheme's
strictness and teaching (the strictness scheme at $q = \QGridSOfsted$).

The prediction model predicts enacted warmth from a report's wording, correlating at $\PredWOfsted$, and enacted teaching at $\PredTOfsted$, and both figures
are essentially unchanged, or higher, when everything demographics could explain is removed.
Strictness is different: no prediction model predicts it well. The best attempt reaches
$\PredSOfsted$, it beats luck only narrowly (\PermPOfstedS{}), and its confidence interval
touches zero ($\CIPredSOfsted$). Every word-count representation tried showed the same weakness.

For strictness the mark scheme does better than the prediction model, although the gap between the two is small enough to
be chance
at this sample size ($z = \StrictDiffZ$, $p = \StrictDiffP$). The numbers alone therefore cannot
separate the two methods on strictness. The three checks set out in the measurement section
can: the retest, the change of language model, and the length check. The scheme's strictness bands pass all three.

The retest asks whether the language model gives the same answer twice. Each scheme was run
twice on the same documents. Across the seven mark schemes, the two runs give the same
band \GridRetestExact{} of the time, and the report strictness scheme gave every school
the same band on both runs.

The purpose of changing the language model is to check whether the reading is genuinely down
to the mark scheme or whether the particular model that applies it makes a difference.
If the scheme is doing the work, a different model applying the same mark scheme should find the same thing. For strictness, the second model's bands correlate with the enacted scores at $\RubricSFourO$, close to the primary model's $\RubricSOfstedRaw$, so the
scheme is what determines the band, not the model. The warmth and teaching schemes fail
the same check. Applying the same mark schemes, the second model's bands correlate with the enacted scores at essentially nothing ($\RubricWFourO$ for warmth, $\RubricTFourO$ for teaching), so what the
primary model found there came from the model rather than the scheme.

The length check asks whether the mark-scheme bands follow document length rather than the school. For the inspection-report schemes the worry is small: the strictness bands correlate with length at $\LenOfstedS$ and the warmth bands at $\LenOfstedW$, and removing the effect of length moves the strictness correlation only from $+0.25$ to $\RubricSOfstedLenP$. The website schemes' bands follow length closely, one more reason their bands are not used.

Demographic adjustment is separate from the three checks, and it is applied to the mark-scheme bands in the same way as to the prediction models' predictions; \cref{sec:p1_app_prediction} gives the specification. The adjustment changes the strictness correlation little: the scheme's correlation with enacted strictness is $\RubricSOfstedRaw$ raw and $\RubricSOfstedDec$ adjusted. The warmth scheme's bands, by contrast, correlate with enacted warmth at $\RubricWOfstedDec$ and put three-quarters of schools in one band. The model, on the same reports, predicts enacted warmth at $\PredWOfsted$,
and its advantage is itself significant (paired test, $p = \WarmDiffP$).

The two dimensions appear in the reports differently. Warmth is associated with many small differences in word choice across the whole report. A model that counts words can learn these from the visited schools, but a mark scheme cannot list them in advance. Strictness is associated with a few specific claims, for example that disruption is rare or that rules are applied consistently. A reader must be told to look for these claims. A word-count model cannot use them, because a report that says disruption is rare and one that says it is not rare use almost the same words. Teaching, kept in the
analysis as the rival explanation for the next chapter, sits in between and is read about
equally well by both methods. However, only the model's predictions pass the change-of-model check (the second model's teaching bands correlate at $\RubricTFourO$), so
the prediction model scores are used for teaching.

Behaviour policies and websites tell neither method anything about enacted culture. The prediction
models reach at most $\OwnVoiceMax$ on any pairing, too little to distinguish from zero at this
sample size, and the bands from the mark schemes correlate less still. One of those two failures needs a caution. The
two schemes that read the behaviour policies, one for warmth and one for strictness, each put
nearly every school in the same band, and a scheme that gives nearly every school the same score
cannot correlate with anything, whatever the policies contain. Its zero is therefore not evidence that the policies say nothing. The evidence comes from the prediction models. They give different schools different scores, so a signal in the policies could have shown up, and none does.

\subsection{Which parts of a report the prediction models use}

The prediction models learn which words and phrases predict the enacted scores, and nothing in the method restricts them to any topic. A model can therefore predict the enacted scores accurately for two reasons. The first is that it has picked up the words inspectors use when they describe the dimension itself, such as the way staff and pupils treat one another in the case of warmth. The second is that it has picked up the positive tone that runs through a report on a good school whatever the topic. Warm schools tend to be good schools, so a model that has learned only the general tone would still predict warmth, but it would not be reading warmth. As an informal check on the method, I annotated the content of the reports and asked whether the sentences that raise each model's predicted score most are consistent with the dimension it is meant to measure. For example, are the sentences that predict warmth about relationships between staff and pupils? For the annotation, a language model read every sentence of every report and tagged what the sentence is about: relationships, behaviour, teaching, or something else. The tagging was produced without seeing the schools' scores or anything the models learned, so it can be used to check what they learned.

The check takes the hundred sentences that raise the warmth model's predicted score most, where a sentence's contribution is the sum of the weights the model has learned for the phrases in it, and asks what they are about. \WarmTopShare{} of them are about staff--pupil relationships, against a base rate of \WarmBaseShare{} across all sentences, so the sentences that raise the warmth score are about relationships far more often than chance would give. For teaching the corresponding shares are \TeachTopShare{} against \TeachBaseShare{}, which is no enrichment at all. Inspection reports are mostly about teaching already, so the sentences that raise the teaching score cannot be about teaching much more often than the rest are. The teaching model therefore picks up the report's general tone rather than any one topic.

The topic shares, together with the demographic adjustment reported earlier, say what a model's score is made of. Part of it is language specific to the model's dimension: for warmth, the relational sentences the model seeks out. The rest is the general tone of a school doing well, which runs through a good school's whole report whatever the sentence is about. The demographic adjustment removes some of that general tone, but only the part that the four control variables (three intake measures and report age) can predict; the remainder stays in the score. A predicted warmth score therefore combines two things. Part is the relational language specific to warmth. Part is the general positive tone of a report on a good school. It is a measure of the report's whole account of the school, weighted towards warmth, not a measure of warmth alone.

\Cref{tab:p1_examples} shows the sentences each model weighs most heavily, under selection rules
stated with the table. The two listings differ. For teaching the strongest sentences name the
topic directly. For warmth they do not.

For teaching, the strongest sentences are about pedagogy: choosing core knowledge, embedding
learning in long-term memory, pupils focusing in lessons. Some of the strongest, though, are
criticisms. A sentence saying that pupils are not remembering knowledge deeply raises the score
as surely as one saying that they are. That happens for two reasons. The first is a flaw in the model. It counts words and cannot tell praise from
criticism. The second is a fact about reports. Inspectors write about long-term memory, whether in praise or criticism, only when a school teaches in that way, so the vocabulary is informative either way. The appendix discusses both in more detail.

For warmth the pattern is different, and at first sight stranger. The sentences that most lower
a school's warmth score sound kind. ``Many pupils feel confident they have a `trusted adult' they
can go to'' and ``pastoral support is in place'' are among the strongest score-lowering sentences.
The explanation for this is that an inspector writes about a warm school by describing adults and
pupils dealing with one another. When there is nothing like that to describe, the inspector
mentions support systems instead. Systems are what gets written about when the relationships
are missing.

On the positive side, no single sentence carries much weight. The model's warmth evidence is
spread across many small choices of phrase: attentive verbs, pupils' names used, care described
in passing. Even its strongest whole sentences read as the tone of an engaged community rather
than tenderness. The spread is a fact about warmth, not a weakness of the model. Warmth never
concentrates in one sentence, so there is no quotable line to point to. The case that the model
is reading warmth therefore rests on two things shown already. The first is the held-out
correlation with the warmth scores for the visited schools. The second is the topic shares:
the sentences that raise the model's predicted warmth most are
about staff--pupil relationships far more often than their share of all sentences would
predict. \Cref{sec:p1_app_examples} also lists the sentences about warmth towards staff that raise the predicted score: schools whose leaders listen to their teachers are predicted to be warmer. That
part of the signal is easy to interpret, and it is not difficult to imagine that a school
where leaders listen to staff is also one where teachers listen to the concerns of their pupils.

% generated by make_ch2_examples.py -- do not edit
\begin{table}[htbp]
\centering
\caption{Sentences from the visited schools' inspection reports that most raise or lower the models' scores. A sentence's score is the sum of its phrases' learned weights; the full listing is in \cref{sec:p1_app_examples}.}
\label{tab:p1_examples}
{\small\begin{tabular}{p{0.93\textwidth}}
\toprule
\multicolumn{1}{l}{\emph{Warmth --- raises the score}} \\
``They are taught how to show respect for different opinions and listen to the views of others.'' \\
``They hold leaders strictly to account for the quality of the school's provision and are highly effective.'' \\
``They listen to others and express their views in a well-managed environment.'' \\
``Sixth-form students are proud role models and make a highly positive contribution to the school and local community.'' \\
\addlinespace
\multicolumn{1}{l}{\emph{Warmth --- lowers the score}} \\
``However, some pupils, particularly some disadvantaged pupils, are absent too often.'' \\
``Many pupils feel confident they have a `trusted adult' they can go to.'' \\
``(Information for the school and appropriate authority)  In some subjects, the work given to some pupils does not enable them to develop their ideas and extend their answers.'' \\
``Pupils take part in a range of extra-curricular activities to develop their talents and interests, for example basketball, choir, photography, and cookery.'' \\
\addlinespace
\multicolumn{1}{l}{\emph{Teaching --- raises the score}} \\
``Those responsible for governance know the school's context and provide support and professional challenge where necessary to ensure that the school's high standards are maintained.'' \\
``Leaders are supported and challenged by dedicated and knowledgeable governors, who share their vision of high standards in a supportive school.'' \\
``In a few subjects, however, pupils are not remembering knowledge deeply in the long term.'' \\
``The school's motto is `excellence with care.' Pupils testify to how this works in practice.'' \\
\addlinespace
\multicolumn{1}{l}{\emph{Teaching --- lowers the score}} \\
``The personal, social and health education (PSHE) curriculum is designed to support pupils' wider development.'' \\
``Pupils take part in a range of extra-curricular activities to develop their talents and interests, for example basketball, choir, photography, and cookery.'' \\
``On other occasions, adults do not support pupils to complete work to a high standard.'' \\
``Judgements in this report are based on the current inspection framework and also reflect changes that have happened at any point since the last graded inspection.'' \\
\addlinespace
\bottomrule
\end{tabular}}
\end{table}


\subsection{What schools write about themselves}

A school's own documents, its behaviour policy and website, tell a reader what the head says about the school, not what the school does. The test is a rerun of the prediction models: they read the policy and website text as before, but this time they are trained to predict the head's interview answers rather than the enacted scores. This was done across the roughly 300 interviewed schools. On strictness the documents match what the head says: the models predict the head's account at $\PredEspS$ (demographics removed), which survives the correction for the number of tests. On staff morale the relation is weaker, at $\PredEspClim$: significant on its own, but not after that correction. On warmth there is no relation at all, at $\PredEspW$. What a school writes about warmth is unrelated even to what its own head says in the interview.

The bands from the website mark schemes do not correlate with any enacted or espoused score, on any dimension. That is weak evidence rather than contrary evidence: two language models applying the same website mark scheme agree on the band only \AgreeWebW{} and \AgreeWebS{} of the time, against \AgreeOfstedW{} to \AgreeOfstedT{} on inspection reports, so the scheme is barely reliable on the documents schools write about themselves.

One caveat qualifies that last result. The espoused warmth score is built from just three interview
statements, and heads bunch their answers near the top of the scale, so the score varies little
from school to school. When the target score varies little between schools, any correlation with it is lower than it would be if the scores were spread out. Finding no correlation against such a target therefore rules out a strong link between public warmth language
and the head's account, but not a modest one.

The inspection report at first seems to break this pattern. Its text correlates with the head's account of strictness at $\CrossWireRaw$, as if the inspector somehow knew what the head thinks. There is a simpler explanation. The report describes what the school does, and what the school does is itself related to what the head says. If that is the route, the report is linked to the head only through the school, and the inspector never needed to know the head's views at all. The test removes from both measures the part that enacted strictness predicts and correlates what remains. The correlation falls from $\CrossWireRaw$ to $\CrossWireNet$. So much of the link between the report and the head's account is explained by the school's actual behaviour, which both describe. The report describes the school. The head's account also describes the school. That is why the two agree, and the inspector did not need to know the head's views. So the two kinds of document line up with the two measures. What a school writes about itself matches what its head says, the espoused side. What the inspectorate writes about the school matches what the school does, the enacted side.

\subsection{A check against the parent survey}
\label{sec:p1_pv}

The last check comes from parents. Everything so far has judged the prediction model's warmth scores against the visited schools, and there are only 103 of them. Parents offer a second, independent judge at national scale: Ofsted's Parent View survey asks them, at nearly every school in England, whether their child is happy at the school and feels safe there, which is warmth as a parent experiences it. If the model's warmth scores are real, they should agree with what parents report. This section makes that comparison and reports two results: the model's warmth scores do agree with the parents, and the agreement is not explained by the inspection grade.

The measure is the parent-reported warmth score: the share of parents who say their child is happy at the school and feels safe there. There are two important limitations to this data that should be kept in mind.
The first is that parents choose whether to
respond, and the typical school's figure rests on about 120 replies, roughly one parent in eight.
The second is a timing issue. Responses are taken over several years, so at some schools the
parents' answers describe an earlier version of the school than the one the culture data describe.

The first result is the national agreement. The warmth prediction model was estimated on the visited schools alone. That leaves \PVPredWarmthN{} schools the model
has never seen. For each of these, the model reads the school's inspection report and turns it
into a predicted warmth score. Those predictions can then be set against what parents say.
Across the \PVPredWarmthN{} schools, predicted warmth and parent-reported warmth correlate at $\PVPredWarmth$; \cref{fig:p1_pv} plots the two measures against each other. That agreement would mean nothing if the model had been given the answer: two measures agree trivially when one was built from the other, as when a model
is tested on the data it was trained on. Neither is the case here. The model never saw these schools when it was estimated, and it never saw any parent data at all. The only way it can agree
 with the parents is if both are picking up something true about the school.

\begin{figure}[htbp]
  \centering
  \caption{Predicted warmth and parent-reported warmth, nationally}
  \label{fig:p1_pv}
  \includegraphics[width=0.55\textwidth]{fig_ch2_pv.pdf}
\begin{minipage}{\linewidth}\smallskip
\footnotesize\textit{Notes:} Correlation $r = \PVPredWarmth$ on $n = {}$\PVPredWarmthN{} schools. Each dot is a bin of schools (twenty equal-count bins); the model never saw any of these schools in training.
\end{minipage}
\end{figure}

The second result concerns the inspection grade. The model reads inspection reports, and a report reflects the grade the school received. If the model's warmth score were the grade in different words, the parents' agreement with the model would add nothing beyond the parents' agreement with the grade. The test is to compare schools that share the same grade. Within such a group the grade is the same for every school, so it cannot explain why the model scores one school warmer than another. On the visited schools, 99 of which have a grade from an inspection before the pandemic, the 2019 grade used in the next chapter, the model's correlation with enacted warmth is $+0.32$, essentially its full-sample value of $+0.33$. Comparing only schools with the same grade leaves it at $+0.29$. The teaching model gives the same result. Its correlation with enacted teaching is essentially unchanged when only schools with the same grade are compared. The model's warmth scores are not the grade in different words. \Cref{sec:p1_app_parents} adds a stricter version of this check, together with what the visited schools show and a check on one possible route: some reports quote the Parent View survey, so the model could in principle be reading the parents' answers
rather than independently agreeing with them.

The wider point is that the two measures, the parents' survey answers and the model's warmth scores, share nothing that could make them agree for any reason other than the school.
There is a final point to consider, though.
A school that is good in general tends to have both a report the model scores as warmer and happier parents, so
part of the agreement may reflect general quality rather than warmth in particular. The
agreement therefore stands as corroboration rather than proof.

\section{Conclusion}
\label{sec:p1_conclusion}

Warmth and strictness can both be measured by observation and in public text. Warmth is the harder
of the two to measure. In person it needs more lessons and more settings than
strictness does to reach the same reliability. In public text it is detected only by a model that counts words and phrases. The words associated with warmth are many and each matters a little, so they are spread across an inspection report's ordinary wording, in whether the inspector
writes about adults attending to pupils or falls back on stock assurances. A mark scheme asks the language model to find sentences that describe warmth, and no single sentence does.

Strictness is the opposite. In person it is easy to observe reliably. In text it sits in specific
claims: whether disruption is rare, whether rules are enforced consistently. Checking claims is necessary
because a report saying that disruption is rare and one saying that it is common use
almost the same words. That is why the prediction model predicts strictness weakly but the mark scheme does better.

The documents schools write about themselves match what heads say about strictness and staff morale. On warmth they carry no signal detectable at this sample size, either against the
observers or against what headteachers say when interviewed.

The parent surveys provide a useful check. At the national level their reports of happy, safe children
agree with the model, which supports the claim that inspection reports carry a real warmth signal.
In the visited schools, as the appendix shows, parents' answers are closer to the head's account than to the observers' scores.

There are four limitations that should be noted. The first is that the visit scores measure
warmth and strictness with error. Two observers watching the same school agree well but not
perfectly, so part of each score is error, and every correlation with the visit scores is lower than it would be without that error. The reported correlations therefore understate the true relations. The second is that inspection reports are written in different years under
different inspection frameworks, which adds noise. Removing dates, boilerplate and framework-specific wording reduces that noise but cannot remove it.
This kind of noise pushes correlations towards zero.

The third limitation is that each school was visited once. A visit is a
snapshot, and a school may have had an unusually calm or an unusually rough day; the scores
cannot tell an unusual day from an ordinary one.

The final limitation is that the model's warmth scores cannot be fully separated from the inspection's overall verdict. The prose
the model reads and the grade the inspection awarded come from the same visit and say
related things, so part of what the model reads may be the inspection's judgement of the
school rather than its warmth in particular. The checks in the results show that the model
reads more than the grade alone. What they cannot show is how much of the rest is warmth,
because the grade and the prose come from the same visit, so the grade already reflects some of the warmth the prose describes. Any comparison that removes the grade therefore removes part of the warmth signal as well.

With those limitations noted, the chapter provides three national measures, each validated against direct observation: prediction models for warmth and teaching, and a one-page mark scheme
for strictness. The following chapter uses these to investigate the associations between
warmth, strictness, teaching and pupil progress.

\startchapterappendix{2}
\begin{center}
  {\LARGE\bfseries Appendix to Chapter 2}\\[6pt]
  \rule{0.5\linewidth}{0.4pt}
\end{center}
\medskip
\noindent This appendix collects the supporting material for the chapter,
in the order in which the chapter uses it. It sets out how the visit and
interview items combine into scores, reproduces the two observation forms and
the interview guide as they were administered, gives the inter-rater
reliability tables, describes the samples and the coverage of the text
collections, specifies the prediction models, lists the sentences and phrases those models weigh most heavily, collects the further checks on the parent comparison, and finally reproduces the text of every mark scheme given to the language model.
\medskip

% - - - - - - - - - - - - - - - - - - - - - - - - - - - - - - -
\section{Scoring Guide and Item Definitions}
\label{app:2A:guide}

This section sets out how the individual items on the forms combine into
the nine sub-scores and the two sets of scores. The items themselves, that
is, the question wording, the response format and the anchor description
attached to each rating level, are reproduced in full in
\cref{app:2A:interview_guide} for the interview and in
\cref{app:2A:visit_protocol} for the two observation forms.

Every item is recorded on a 1--5 scale. Two observation items can also take
a zero, with a distinct meaning in each case. A zero on \emph{response of
pupils to sanctions} records that no sanction was issued during the lesson,
and a zero on \emph{effectiveness of pupil discussion} records that no
discussion was initiated. For the reasons given in the chapter, the first of these is not a genuine low score. It is replaced by $6 - \text{disruption}$, which gives a calm unsanctioned lesson a high value
and a disrupted unsanctioned lesson a low one. The second is treated as a
genuine negative, since a lesson containing no discussion did not achieve
effective discussion.

Observation item scores are averaged across lessons and across researchers
to give one value per school before the sub-scores are formed.
\Cref{tab:p1_subscore_items} lists the constituent items of each sub-score.

% Auto-generated by thesis/make_instruments.py --- do not edit by hand.
% Source: warm_strict_scorer.py

\begin{longtable}{@{}p{0.055\linewidth}p{0.235\linewidth}p{0.62\linewidth}@{}}
\caption{Composition of the gold-standard sub-scores. Item wording is given in full in \cref{app:2A:interview_guide} and \cref{app:2A:visit_protocol}.}
\label{tab:p1_subscore_items}\\
\toprule
& Sub-score & Constituent items \\
\midrule
\endfirsthead
\toprule
& Sub-score & Constituent items \\
\midrule
\endhead
\bottomrule
\endfoot
\multicolumn{3}{@{}l}{\textbf{Warmth}} \\
W1 & Teaching warmth & \textit{Lesson observation.} Frequency with which teacher uses pupils' names; Frequency of praise; Quality of interactions between teachers and pupils \\
W2 & School warmth & \textit{Outside-of-lessons observation.} Quality of relationships between pupils during transition; Quality of non-behaviour related interactions between staff and pupils at break; Quality of relationships between pupils at break \\
W3 & Headteacher warmth philosophy & \textit{Interview.} Q7 Extracurricular provision; Q11 Pupil wellbeing; S2 Reward systems applied consistently; S7 Pupils feel safe at school; S8 Pupils feel their teachers care \\
\addlinespace
\multicolumn{3}{@{}l}{\textbf{Strictness}} \\
S1 & In-lesson strictness & \textit{Lesson observation.} Level of misbehaviour tolerated before a first sanction (reversed); Frequency of low-level disruption (reversed); Response of pupils to sanctions; Extent to which pupils address the teacher respectfully \\
S2 & Out-of-lesson strictness & \textit{Outside-of-lessons observation.} Consistency of use of sanction system; Calmness of corridors; Arrival time to next lesson; Organisation of canteen/dining hall; Organisation of recreational space \\
S3 & Systems strictness & \textit{Interview.} Centralised detentions; Internal exclusion room; Behaviour points/demerits; Line-ups; Silent corridors; After-school revision; Weekend revision; Ban on mobile phones \\
S4 & Headteacher strictness philosophy & \textit{Interview.} Q5 Reward and sanction system; S1 Sanction systems applied consistently; S5 Pupil behaviour in class is generally good; S6 Pupil behaviour outside class is generally good; S9 Teachers feel respected by pupils \\
\addlinespace
\multicolumn{3}{@{}l}{\textbf{Teaching practice}} \\
T1 & Teaching practice (observed) & \textit{Lesson observation.} Frequency of questioning/checking pupil progress; Quality of verbal feedback; Effectiveness of pupil discussion; Quality of differentiation; Clarity of teacher explanation; Extent to which learning outcomes are achieved; Effectiveness of teaching methods; Clarity of lesson structure; Effectiveness of resource use \\
T2 & Teaching practice (stated) & \textit{Interview.} Q6 Characteristics of an outstanding teacher; Q8 Marking and feedback; Q12 Assessing pupil progress; S3 Teacher turnover is low; S4 Staff morale is good \\
\end{longtable}


Every other sub-score is the mean of its items, computed over the items
present so that a missing item reduces the number of contributing items
rather than the school's score.

The sub-scores combine into two sets of scores, grouped by source, with
observed behaviour and stated philosophy kept apart:
%
\begin{align}
\text{Enacted:}\quad
  W^{\text{enac}} &= \overline{(W1, W2)}, &
  S^{\text{enac}} &= \overline{(S1, S2)}, &
  T^{\text{enac}} &= T1 \label{eq:enacted} \\
\text{Espoused:}\quad
  W^{\text{esp}}  &= W3, &
  S^{\text{esp}}  &= S4, &
  \mathrm{SC} &= \overline{(\text{SC}_{1}, \text{SC}_{2})} \label{eq:espoused}
\end{align}
%
There is no espoused teaching score. The two interview statements under the
teaching heading concern staff turnover and morale, and they form the staff-climate score SC instead. The systems count S3 enters no espoused score.

Schools with an interview but no visit have the espoused set of scores alone
(every visited school was also interviewed). No imputation is attempted across the two sets: as the chapter reports, the espoused and enacted scores are too weakly correlated for either to stand in for the other.

The mapping in \cref{tab:p1_subscore_items} is generated directly from the
scoring code, so it cannot drift from what the scoring code does.

The assignment of items to dimensions is theoretical, but it can be checked
against the data. A two-factor analysis of the school-level item means
(principal components with varimax rotation, on the 98 schools with every
culture item observed) recovers most of the structure Table~2A.1 imposes:
fourteen of the nineteen culture items load most heavily on their assigned
dimension, with the in-lesson strictness items and the relational items
separating most cleanly. The remaining five load more on the other dimension. Those five are either near-ties or weakly loaded items about order in transitions. When the ten teaching items are
added they form the dominant factor, and eighteen of the nineteen culture
items then sort with their assignment. The analysis forces the two factors to be uncorrelated, while the chapter estimates the two dimensions correlate at about $0.36$. Items that belong to both are therefore harder for this method to assign, and the check is if anything conservative.

\clearpage
% - - - - - - - - - - - - - - - - - - - - - - - - - - - - - - -
\section{School Visit Protocol}
\label{app:2A:visit_protocol}

% Auto-generated by thesis/make_instruments.py --- do not edit by hand.
% Source: Novel Data/School Visit Proforma - *.docx

Two proformas were used on each visit. The lesson observation form was
completed once per lesson observed, and the outside-of-lessons form once per
school, covering transitions, break and lunch. Every item is rated 1--5 against
the anchor descriptions below, with anchors specified at 1, 3 and 5 and even
values used for intermediate judgements. Visits were made by teams of two or
three researchers, who completed the forms independently; item scores are
averaged across researchers and across lessons to give one value per school,
and \cref{sec:p1_irr_obs} reports agreement between observers.

\subsection*{Lesson observation}

\subsubsection*{Section 1 --- Strictness and Warmth}

\paragraph{Level of misbehaviour that results in a first sanction/level of misbehaviour which is tolerated}
Rated 1 (minor) to 5 (serious).

\begin{description}
\item[1] Isolated instances of misbehaviour such as whispering or fidgeting.
\item[3] Repeated disruptions, e.g., repeated talking or distraction of peers.
\item[5] Aggressive or defiant behaviour
\end{description}

\paragraph{Response of pupils to sanctions}
Rated 1 (very poor) to 5 (very good). A score of 0 is recorded separately (see below).

\begin{description}
\item[0] If no sanctions are observed
\item[1] Students ignore sanctions or react defiantly, showing no change in behaviour.
\item[3] Noticeable change in behaviour in response to sanctions, but not consistently maintained.
\item[5] Immediate and sustained positive response to sanctions e.g. apologises for actions.
\end{description}

\paragraph{Frequency of low-level disruption}
Rated 1 (very low) to 5 (very high).

\begin{description}
\item[1] Classroom is calm with disruptions rare or not occurring at all.
\item[3] Regular disruptions, sometimes affecting the flow of the lesson but generally manageable.
\item[5] Constant disruptions, creating a challenging environment for teaching and learning.
\end{description}

\paragraph{Concentration level of students}
Rated 1 (very low) to 5 (very high).

\begin{description}
\item[1] Students are consistently off-task, showing little interest or engagement with the lesson.
\item[3] Moderate levels of concentration, with students generally engaged but susceptible to distractions.
\item[5] Exceptional concentration, with students deeply engaged and consistently focused throughout the lesson.
\end{description}

\paragraph{Frequency of teacher-initiated participation}
Rated 1 (very low) to 5 (very high).

\begin{description}
\item[1] Teacher rarely encourages student participation, leading the class in a predominantly lecture-style format.
\item[3] Regular attempts by the teacher to involve students in the lesson, but some parts where teacher talks continuously.
\item[5] Teacher consistently involves students in the learning process.
\end{description}

\paragraph{Frequency of student-initiated participation}
Rated 1 (very low) to 5 (very high).

\begin{description}
\item[1] Students rarely or never take the initiative to participate, showing passivity in the learning process.
\item[3] Moderate frequency of student-initiated participation. Sometimes active but sometimes passive.
\item[5] Students consistently participate in the lesson demonstrating high levels of engagement in the learning process.
\end{description}

\paragraph{Motivation level of students}
Rated 1 (very low) to 5 (very high).

\begin{description}
\item[1] Students exhibit a lack of interest and motivation, showing little effort in classroom activities.
\item[3] A moderate level of motivation, with students showing interest at times but being disengaged at times.
\item[5] Exceptional motivation and enthusiasm, with students consistently demonstrating a strong desire to learn and participate.
\end{description}

\paragraph{Extent to which students address the teacher respectfully}
Rated 1 (not at all) to 5 (always).

\begin{description}
\item[1] Disrespectful interactions, use of inappropriate language or tone.
\item[3] Respectful interactions most of the time
\item[5] Exceptionally respectful and considerate interactions at all times.
\end{description}

\paragraph{Frequency with which teacher uses pupils' names}
Rated 1 (very low) to 5 (very high).

\begin{description}
\item[1] Rarely addresses students by name; impersonal interactions.
\item[3] Regularly uses student names, but some students are less frequently addressed.
\item[5] Consistently uses every student's name, fostering a sense of individual attention.
\end{description}

\paragraph{Frequency of praise}
Rated 1 (very low) to 5 (very high).

\begin{description}
\item[1] Teacher provides little to no praise, no positive reinforcement.
\item[3] Regular praise, but it is not always meaningful or specific.
\item[5] Regular and specific praise which highlights what students have done well.
\end{description}

\paragraph{Quality of interactions between teachers and pupils}
Rated 1 (impersonal) to 5 (warm).

\begin{description}
\item[1] Interactions are formal, distant, and lack warmth.
\item[3] Interactions are warm, but still quite formal in nature.
\item[5] Exceptionally warm and supportive interactions throughout the lesson.
\end{description}

\subsubsection*{Section 2 --- Teaching practice}

\paragraph{Frequency of questioning/checking pupil progress}
Rated 1 (very low) to 5 (very high).

\begin{description}
\item[1] Teacher rarely asks questions or checks on pupils' understanding.
\item[3] Teacher regularly asks standard questions, adequately assessing pupils' basic understanding but may not probe deeper comprehension.
\item[5] Teacher consistently integrates questioning seamlessly into lessons, skilfully using them to assess understanding and encourage higher-order thinking.
\end{description}

\paragraph{Quality of verbal feedback}
Rated 1 (very poor) to 5 (very good).

\begin{description}
\item[1] Feedback is rare, nonspecific, or irrelevant, offering no guidance or improvement for pupil learning.
\item[3] Feedback is regularly given, is clear, and somewhat constructive.
\item[5] Feedback is exceptionally well-tailored to each pupil and highly constructive.
\end{description}

\paragraph{Effectiveness of pupil discussion}
Rated 1 (not effective) to 5 (very effective). A score of 0 is recorded separately (see below).

\begin{description}
\item[0] No teacher-initiated pupil discussion is observed.
\item[1] Pupil discussions are rarely initiated or if they are, then they are poorly managed, resulting in minimal pupil engagement.
\item[3] Regularly encourages discussions which have a moderate impact on learning but may not fully engage all pupils.
\item[5] Regularly encourages discussions and structures them effectively to ensure that pupils engage with them and benefit from them.
\end{description}

\paragraph{Quality of differentiation (adaptation to students' needs)}
Rated 1 (very poor) to 5 (very good).

\begin{description}
\item[1] Teaching lacks adaptation to different learning needs, often leading to disengagement or unmet learning objectives for various students.
\item[3] Regularly employs differentiation strategies that address the needs of most students but not all.
\item[5] Exceptional in adapting teaching to meet all students' needs. Challenges high ability pupils and/or supports low ability pupils appropriately.
\end{description}

\paragraph{Clarity of teacher explanation}
Rated 1 (very poor) to 5 (very good).

\begin{description}
\item[1] Explanations are often unclear, leading to widespread confusion or misconceptions among students.
\item[3] Usually provides explanations that are clear and aid in student understanding, though some concepts may not be understood by all students.
\item[5] Provides exceptionally clear and comprehensive explanations, making complex concepts accessible to all or almost all students.
\end{description}

\paragraph{Extent to which learning outcomes are achieved}
Rated 1 (not at all) to 5 (completely).

\begin{description}
\item[1] Learning outcomes are rarely met, with most students failing to achieve the set objectives. Unclear what learning objectives are.
\item[3] A majority of learning outcomes are met, though some objectives may not be fully achieved by all students.
\item[5] All learning outcomes are fully met, with evidence of deep understanding and application of knowledge by all or almost all students.
\end{description}

\paragraph{Effectiveness of teaching methods}
Rated 1 (not effective at all) to 5 (very effective).

\begin{description}
\item[1] Teaching method(s) are not appropriate to the content and tasks.
\item[3] Teaching method(s) are mostly appropriate to the content and tasks.
\item[5] Chosen teaching method(s) are very well matched to the content and every task is clearly well thought out.
\end{description}

\paragraph{Clarity of lesson structure}
Rated 1 (unclear) to 5 (very clear).

\begin{description}
\item[1] Lessons lack a clear structure, resulting in disorganized content delivery. Structure is not communicated to students.
\item[3] Generally clear lesson structure with a logical flow, though minor issues in transition or clarity may occasionally arise. Structure is communicated quite well to students.
\item[5] Lessons are exceptionally well-structured, with seamless transitions, and a well-thought-out progression. Structure is regularly communicated to students.
\end{description}

\paragraph{Effectiveness of resource use}
Rated 1 (low) to 5 (high).

\begin{description}
\item[1] The resources do not allow most students to achieve the learning objectives and/or engage with the topic.
\item[3] The resources allow most students to achieve the learning objectives and/or engage with the topic.
\item[5] The resources allow all or almost all of the students to achieve the learning objectives and/or engage with the topic.
\end{description}


\subsection*{Outside-of-lessons observation}

\subsubsection*{Section 1 --- Comparison across lessons}

\paragraph{Consistency of use of sanction system}
Rated 1 (not consistent at all) to 5 (completely consistent).

\begin{description}
\item[1] Teachers are using different sanction systems in every classroom.
\item[3] Teachers are generally taking the same approach to misbehaviour, but there is no shared language.
\item[5] All teachers use the sanction system in exactly the same way and there is a clear shared language.
\end{description}

\paragraph{Consistency of use of reward system}
Rated 1 (not consistent at all) to 5 (completely consistent).

\begin{description}
\item[1] Teachers are using different reward systems in every classroom.
\item[3] Teachers are generally taking the same approach to positive behaviours, but there is no shared language.
\item[5] All teachers use the reward system in exactly the same way and there is a clear shared language.
\end{description}

\paragraph{Consistency of verbal feedback}
Rated 1 (not consistent at all) to 5 (completely consistent).

Should be calculated as the average of the scores from individual lessons.

\paragraph{Quality of differentiation (adaptation to students' needs) across lessons}
Rated 1 (very poor) to 5 (very good).

Should be calculated as the average of the scores from individual lessons.

\subsubsection*{Section 2 --- Transitions}

\paragraph{Calmness of corridors}
Rated 1 (chaotic) to 5 (calm).

\begin{description}
\item[1] The corridors feel unsafe. There is a lot of boisterous behaviour and students do not seem concerned with getting to their next lesson.
\item[3] Most students move between lessons in a calm fashion. There is an occasional loud voice and a small amount of boisterous behaviour. Most students move purposefully to their next lesson.
\item[5] All students move between lessons in a calm fashion with no loud voices or boisterous behaviour at all. All students move with purpose to their next lesson.
\end{description}

\paragraph{Arrival time to next lesson}
Rated 1 (lots of lateness) to 5 (all on time).

\begin{description}
\item[1] The majority of students arrive late to lessons.
\item[3] Most students arrive on time with a few late arrivals.
\item[5] All or almost all students arrive on time.
\end{description}

\paragraph{Quality of relationships between students during transition}
Rated 1 (very poor) to 5 (very good).

\begin{description}
\item[1] Lots of students speak to each other disrespectfully. There is clear evidence of bullying and/or intimidation.
\item[3] Most students are polite in their interactions with each other. There is evidence that a small amount of bullying/intimidation is taking place.
\item[5] All or almost all students are polite in their interactions with each other. There is no evidence of bullying or intimidation of any kind.
\end{description}

\paragraph{Presentation of wall/corridor displays}
Rated 1 (poor condition) to 5 (excellent condition).

\begin{description}
\item[1] There are no displays, or if there are, then they are old and tattered.
\item[3] Most displays are in good condition, but some are a bit tattered and look old.
\item[5] All displays are in pristine condition with no damage. They all look as if they are new.
\end{description}

\paragraph{Extent to which wall/corridor displays which are not student work align with school values}
Rated 1 (random) to 5 (complete alignment).

\begin{description}
\item[1] There are no displays which convey messages about the school values.
\item[3] There are some displays which convey messages about the school values.
\item[5] All or almost all displays convey messages about the school values.
\end{description}

\subsubsection*{Section 3 --- Break and lunch}

\paragraph{Quality of non-behaviour related interactions between staff and students}
Rated 1 (very poor) to 5 (very good).

\begin{description}
\item[1] In general, staff do not engage with students.
\item[3] Most staff actively engage with students and interactions are mostly warm and personal.
\item[5] All or almost all staff actively engage with students and interactions are warm and personal.
\end{description}

\paragraph{Quality of relationships between students}
Rated 1 (very poor) to 5 (very good).

\begin{description}
\item[1] Lots of students speak to each other disrespectfully. There is clear evidence of bullying and/or intimidation.
\item[3] Most students are polite in their interactions with each other. There is evidence that a small amount of bullying/intimidation is taking place.
\item[5] All or almost all students are polite in their interactions with each other. There is no evidence of bullying or intimidation of any kind.
\end{description}

\paragraph{Organisation of canteen/dining hall}
Rated 1 (chaotic) to 5 (calm).

\begin{description}
\item[1] Students get their food in a haphazard fashion. They leave food behind on tables and on the floor. They do not take care of their environment.
\item[3] There is some organisation to the way that students get their food. There is a little bit of mess left behind.
\item[5] Students get their food in an organised fashion. They leave the table and floor where they eat clean and tidy.
\end{description}

\paragraph{Organisation of recreational space}
Rated 1 (chaotic) to 5 (calm).

\begin{description}
\item[1] Activities are interfering with each other and there is no supervision from staff. There is a general sense of disorganisation. There is lots of rule breaking.
\item[3] There are clear areas for different activities, and some are supervised by staff. There is some rule breaking.
\item[5] There are clear areas for different activities, and these are all supervised by members of staff. All or almost all students abide by the rules.
\end{description}


\clearpage
% - - - - - - - - - - - - - - - - - - - - - - - - - - - - - - -
\section{Interview Guide}
\label{app:2A:interview_guide}

% Auto-generated by thesis/make_instruments.py --- do not edit by hand.
% Source: Novel Data/Headteacher Interview.docx

The schedule below is the instrument as administered. Each open-ended
question was scored 1--5 against the anchor descriptions shown. Each
question was asked by one researcher and scored by both; \cref{tab:irr_interview}
reports agreement between them, and the analysis uses the final recorded score
for each question. A checklist of key words and phrases was also completed for
each question; it enters no score reported in the chapter and is omitted here.
Item-to-sub-score mapping is given in \cref{app:2A:guide}.

\subsubsection*{Open-ended questions}

\paragraph{Q1. How long have you been in post?}
Recorded verbatim; not scored on the 1--5 scale.

\paragraph{Q2. How would you describe your school's culture?}
Rated 1 (weak answer) to 5 (strong answer).

\begin{description}
\item[1] A vague or generic description is provided, lacking specific details or examples.
\item[3] A detailed description of the school's culture, including some specific examples, but no explanation of the rationale for this culture.
\item[5] A detailed description of the school's culture with specific examples that also includes a clear explanation of the rationale for this culture.
\end{description}

\paragraph{Q3. How do you communicate this culture to your staff?}
Rated 1 (weak answer) to 5 (strong answer).

\begin{description}
\item[1] The strategies used to communicate the culture to the staff are not clearly explained.
\item[3] The strategies used to communicate the culture to the staff are clearly explained, but there is a lack of detail.
\item[5] There is a detailed and clear explanation of the strategies used to communicate the culture to the staff.
\end{description}

\paragraph{Q4. What do you think is the most important characteristic of an effective school? Why?}
Rated 1 (weak answer) to 5 (strong answer).

\begin{description}
\item[1] Various characteristics are mentioned but there is no clear explanation of why these characteristics lead to a school being effective.
\item[3] One or multiple characteristics are mentioned and there is a clear explanation of why these characteristics lead to a school being effective.
\item[5] One or multiple characteristics are mentioned and there is a clear explanation of why these characteristics lead to a school being effective. There is also a discussion of how the school culture links to these characteristics.
\end{description}

\paragraph{Q5. How does your reward and sanction system work? If a pupil misbehaves what should happen? If a pupil makes a positive contribution what should happen?}
Rated 1 (weak answer) to 5 (strong answer).

\begin{description}
\item[1] The explanation of the reward and sanction system is unclear, lacking specific details or a coherent structure.
\item[3] The response outlines the reward and sanction system to a decent extent, but there may be gaps in the explanation, and details might be limited.
\item[5] A clear and detailed explanation is provided on how the reward and sanction system operates; the response demonstrates a deep understanding of the procedures for addressing both misbehaviour and positive contributions.
\end{description}

\paragraph{Q6. What do you think are the key characteristics of an outstanding classroom teacher?}
Rated 1 (weak answer) to 5 (strong answer).

\begin{description}
\item[1] Various characteristics are mentioned but there is no clear explanation of why these characteristics make a teacher outstanding.
\item[3] Various characteristics are mentioned and there is a clear explanation of why these characteristics make a teacher is outstanding.
\item[5] Various characteristics are mentioned and there is a clear explanation of why these characteristics make a teacher is outstanding. There is also a discussion of how these characteristics link to the school's culture.
\end{description}

\paragraph{Q7. How important do you think extracurricular provision is? What kind of extracurricular provision have you put in place?}
Rated 1 (weak answer) to 5 (strong answer).

\begin{description}
\item[1] The importance of extracurricular provision is not clearly articulated, and the response lacks specific details about the types of activities available.
\item[3] The response acknowledges the importance of extracurricular provision, but the explanation may be somewhat basic or lacking in depth; details about the types of activities implemented could be more specific.
\item[5] The response shows a clear understanding of the importance of extracurricular provision and provides specific details about the types of activities in place.
\end{description}

\paragraph{Q8. What is your view on the importance of marking and feedback? Is there a school-wide policy on frequency and type?}
Rated 1 (weak answer) to 5 (strong answer).

\begin{description}
\item[1] The explanation of the importance of marking and feedback is unclear. If there is a school-wide policy, then it is also not explained clearly.
\item[3] The explanation of the importance of marking and feedback is clear. If there is a school-wide policy, then it is also explained clearly but without specific details.
\item[5] The answer lays out a clear and comprehensive view on the importance of marking and feedback. It provides specific details about a school-wide policy, including the frequency and type of marking and feedback.
\end{description}

\paragraph{Q9. What is the structure of your staff team?}
Rated 1 (weak answer) to 5 (strong answer).

\begin{description}
\item[1] The explanation is unclear and difficult to follow.
\item[3] There is a clear explanation of the staff structure from Senior Team through to Middle Leaders.
\item[5] The answer clearly lays out the staff structure from Senior Team down to classroom teachers and support staff with some explanation of the rationale for the structure.
\end{description}

\paragraph{Q10. How do you hold staff to account?}
Rated 1 (weak answer) to 5 (strong answer).

\begin{description}
\item[1] A brief answer which mentions the use of performance management structures but does not go into detail.
\item[3] A clear explanation of the performance management process with discussion of informal and formal steps.
\item[5] A clear explanation of the performance management process with discussion of informal and formal steps, as well as an explanation of an underlying philosophy.
\end{description}

\paragraph{Q11. How do you manage the wellbeing of your pupils?}
Rated 1 (weak answer) to 5 (strong answer).

\begin{description}
\item[1] An imprecise answer which emphasises the importance of managing pupil wellbeing but does not lay out any clear strategies for doing so.
\item[3] An answer which lays out strategies for managing pupil wellbeing, but some strategies are not very precise.
\item[5] An answer which lays out clear strategies for managing pupil wellbeing alongside an explanation of why this is important.
\end{description}

\paragraph{Q12. How do you assess pupil progress?}
Rated 1 (weak answer) to 5 (strong answer).

\begin{description}
\item[1] An answer which mentions no specific strategies for monitoring the progress that pupils are making.
\item[3] An answer which outlines strategies for assessing pupil progress in a summative manner only.
\item[5] An answer which outlines a range of strategies for assessing pupil progress formatively and summatively.
\end{description}

\subsubsection*{Yes-or-no questions}

Does your school make use of the following systems:

\begin{itemize}
\item Centralised detentions
\item An internal exclusion room
\item Reward assemblies/prize-giving ceremonies
\item Behaviour points/demerits
\item Reward points/merits
\item Line-ups
\item Silent corridors
\item After-school revision
\item Weekend revision
\item Ban on mobile phones in school
\end{itemize}

Would you like to say anything about your school's use of any of these systems?

\subsubsection*{Statement ranking}

For each of the following statements circle the appropriate number:

Each statement is rated 1 (strongly disagree) to 5 (strongly agree).

\begin{itemize}
\item The sanction systems are applied consistently by all staff
\item The reward systems are applied consistently by all staff
\item Teacher turnover is low
\item Staff morale is good
\item Pupil behaviour in class is generally good
\item Pupil behaviour outside of class is generally good
\item Pupils feel safe at school
\item Pupils feel like their teachers care
\item Teachers feel respected by all students
\item In general, parents are engaged with their child's education
\end{itemize}

Would you like to say anything about any of your responses to these questions?

Would you be willing to allow a few members of our team to come and observe the normal functioning of your school for a day?

\subsubsection*{To be completed after the interview}

\paragraph{Rate the confidence with which the headteacher responded to the questions}
Rated 1 (not confident at all) to 5 (confident).

\begin{description}
\item[1] There was a lot of hesitation throughout the interview and a lack of clarity in almost all the responses.
\item[3] The answers were generally clear and delivered without hesitation.
\item[5] Every answer was clear and delivered without hesitation. There were no wasted words.
\end{description}

\paragraph{Willingness to reveal information}
Rated 1 (very reluctant) to 5 (totally willing).

\begin{description}
\item[1] There was a reluctance to go into specific details throughout the interview.
\item[3] The headteacher went into specific detail with most answers, but not all.
\item[5] The headteacher was happy to go into very specific detail in every answer.
\end{description}

\paragraph{Interviewee patience}
Rated 1 (little patience) to 5 (lots of patience).

\begin{description}
\item[1] The headteacher rushed through all their answers.
\item[3] The headteacher gave detailed responses to most questions but did not add any detail to the yes/no or statement questions.
\item[5] The headteacher was happy to give detailed responses to all questions and wanted to add detail to their answers to the yes/no and statement questions.
\end{description}


\clearpage
% - - - - - - - - - - - - - - - - - - - - - - - - - - - - - - -
\section{Inter-rater Reliability Tables}
\label{app:2A:irr}

\begin{table}[htbp]
\centering
\caption{Inter-rater reliability for headteacher interview items Q2--Q12.
For each item, a questioner and an observer scored independently on a 1--5 scale.
$r$: Pearson correlation; MAD: mean absolute difference; Exact: proportion of pairs
with identical scores; W1: proportion agreeing within one scale point;
$\kappa_w$: weighted Cohen's $\kappa$ (linear weights).}
\label{tab:irr_interview}
\small
\begin{tabular}{p{5.5cm}rrrrrr}
\toprule
Item & $N$ & $r$ & MAD & Exact\% & W1\% & $\kappa_w$ \\
\midrule
Q2\enspace\textit{{What words describe your school's culture?}} & 292 & 0.709 & 0.56 & 40 & 90 & 0.537 \\
Q3\enspace\textit{{How do you communicate culture to staff?}} & 293 & 0.698 & 0.56 & 33 & 94 & 0.499 \\
Q4\enspace\textit{{What makes an outstanding teacher here?}} & 293 & 0.712 & 0.51 & 41 & 93 & 0.533 \\
Q5\enspace\textit{{How do you reward and sanction pupils?}} & 291 & 0.726 & 0.43 & 44 & 97 & 0.512 \\
Q6\enspace\textit{{What does excellent teaching look like?}} & 293 & 0.613 & 0.52 & 40 & 93 & 0.424 \\
Q7\enspace\textit{{What is your co-curricular offer?}} & 294 & 0.617 & 0.50 & 41 & 95 & 0.448 \\
Q8\enspace\textit{{How do you approach marking and feedback?}} & 294 & 0.691 & 0.48 & 42 & 94 & 0.552 \\
Q9\enspace\textit{{Describe your senior leadership team}} & 294 & 0.623 & 0.45 & 48 & 94 & 0.491 \\
Q10\enspace\textit{{How do you hold staff to account?}} & 291 & 0.707 & 0.52 & 37 & 92 & 0.547 \\
Q11\enspace\textit{{What are biggest challenges facing pupils?}} & 292 & 0.681 & 0.45 & 45 & 93 & 0.550 \\
Q12\enspace\textit{{How do you know how pupils are progressing?}} & 293 & 0.740 & 0.40 & 53 & 95 & 0.606 \\
\bottomrule
\end{tabular}
\end{table}

\begin{table}[htbp]
\centering
\small
\begin{tabular}{lrrrrrr}
\toprule
Item & $N$ & $r$ & MAD & Exact\% & W1\% & $\kappa_w$ \\
\midrule
Misbehaviour & 439 & 0.702 & 0.45 & 51 & 93 & 0.566 \\
Response & 277 & 0.669 & 0.59 & 43 & 88 & 0.545 \\
Disruption & 466 & 0.788 & 0.42 & 53 & 93 & 0.638 \\
Concentration & 467 & 0.647 & 0.51 & 39 & 94 & 0.474 \\
Teacher & 465 & 0.636 & 0.58 & 41 & 89 & 0.450 \\
Student & 457 & 0.637 & 0.63 & 37 & 88 & 0.456 \\
Motivation & 467 & 0.610 & 0.53 & 40 & 93 & 0.454 \\
Respectfully & 467 & 0.654 & 0.53 & 44 & 94 & 0.473 \\
Names & 463 & 0.698 & 0.62 & 43 & 88 & 0.524 \\
Praise & 456 & 0.688 & 0.59 & 39 & 90 & 0.525 \\
Interactions & 463 & 0.616 & 0.62 & 35 & 88 & 0.420 \\
Questioning & 468 & 0.585 & 0.61 & 38 & 89 & 0.440 \\
Verbal & 454 & 0.562 & 0.68 & 37 & 85 & 0.417 \\
Discussion & 217 & 0.722 & 0.59 & 41 & 91 & 0.542 \\
Differentiation & 437 & 0.637 & 0.65 & 38 & 84 & 0.476 \\
Explanation & 454 & 0.578 & 0.51 & 44 & 93 & 0.439 \\
Outcomes & 432 & 0.620 & 0.52 & 44 & 91 & 0.464 \\
Methods & 467 & 0.552 & 0.54 & 41 & 92 & 0.424 \\
Structure & 461 & 0.560 & 0.58 & 40 & 90 & 0.398 \\
Resource & 463 & 0.528 & 0.54 & 43 & 91 & 0.378 \\
\bottomrule
\end{tabular}
\caption{Inter-rater reliability: classroom observation items}
\label{tab:irr_classroom}
\begin{minipage}{\linewidth}
\smallskip
\footnotesize\textit{Notes:} Each row pools all pairwise comparisons across lessons observed by two or three researchers simultaneously. $r$: Pearson correlation; MAD: mean absolute difference; Exact: exact agreement; W1: within-one-point agreement; $\kappa_w$: weighted Cohen's $\kappa$ (linear weights, categories 1--5).
\end{minipage}
\end{table}


\begin{table}[htbp]
\centering
\caption{Inter-rater reliability for outside-of-lesson observation items.
Each row pools all pairwise comparisons across outside-of-lesson observation
sessions attended by two or three researchers. $r$, MAD, Exact, W1, $\kappa_w$
as defined in \cref{tab:irr_classroom}.}
\label{tab:irr_outside}
\small
\begin{tabular}{lrrrrrr}
\toprule
Item & $N$ & $r$ & MAD & Exact\% & W1\% & $\kappa_w$ \\
\midrule
Sanction & 87 & 0.714 & 0.48 & 52 & 91 & 0.610 \\
Reward & 83 & 0.805 & 0.40 & 63 & 90 & 0.706 \\
Verbal & 76 & 0.485 & 0.45 & 32 & 91 & 0.380 \\
Differentiation & 70 & 0.552 & 0.48 & 27 & 93 & 0.345 \\
Corridors & 93 & 0.657 & 0.49 & 39 & 96 & 0.515 \\
Arrival & 91 & 0.474 & 0.58 & 42 & 91 & 0.303 \\
Students & 93 & 0.705 & 0.42 & 47 & 100 & 0.491 \\
Displays & 91 & 0.548 & 0.51 & 54 & 88 & 0.522 \\
Alignment & 84 & 0.578 & 0.62 & 44 & 85 & 0.411 \\
Interactions & 86 & 0.684 & 0.58 & 36 & 92 & 0.511 \\
Relationships & 89 & 0.724 & 0.38 & 54 & 96 & 0.599 \\
Canteen & 91 & 0.727 & 0.38 & 51 & 99 & 0.515 \\
Recreational & 91 & 0.532 & 0.49 & 45 & 92 & 0.446 \\
\bottomrule
\end{tabular}
\end{table}

\clearpage
% - - - - - - - - - - - - - - - - - - - - - - - - - - - - - - -
\section{Sample Characteristics}
\label{app:2A:sample}

\subsection{Sample representativeness}
\label{sec:p1_representativeness}

\Cref{tab:2A:representativeness} and \cref{fig:2A:representativeness_region} compare the
interviewed and visited schools (300 and 102 of them matched to the school panel; 303 and
103 in all) to the national population of 3,312 open secondary and all-through schools in
England, across a range of structural characteristics drawn from the school panel spine and
the Get Information About Schools register, together with the inspection grade as at August
2024 and the pre-pandemic Progress~8 of 2018--19.

The interviewed and visited samples are broadly representative of the national population
on most structural dimensions. School type (academy, maintained, free school), phase of
education, gender composition, admissions policy, and urban--rural classification are all
closely matched to national proportions, with no statistically significant differences.
Mean school size and free school meals percentage are also similar to national averages:
interviewed schools average 1,109 pupils (national mean: 1,079) and a free school meals
share of 27.9\% (national mean: 29.9\%).

The main exception is regional distribution
(\cref{fig:2A:representativeness_region}). Both samples over-represent London and, for
interviews, the South East: taken together, London and the South East account for
44.7\% of interviewed schools versus 30.3\% nationally. The visited sample is
particularly concentrated geographically---London accounts for 45.1\% of visits compared
to 15.0\% of schools nationally---reflecting the logistical constraints of conducting
full-day school visits. Northern England is correspondingly under-represented: the North
West (13.9\% nationally, 8.7\% of interviewed, 2.9\% of visited), the West Midlands
(11.9\%, 7.7\%, 1.0\%), and Yorkshire and the Humber (9.6\%, 5.3\%, 0.0\%) all feature
substantially less in the sample than in the population. Visited schools are also more
likely to have a sixth form (nan\% versus nan\% nationally), consistent with the London
concentration given that London schools disproportionately offer post-16 provision.

Both samples also lean toward stronger schools: 25 per cent of visited schools held an
Outstanding grade as at August 2024 against 15 per cent nationally, and their mean
2018--19 Progress~8 was $+$0.17 against $+$0.04 nationally; the interviewed sample shows the
same lean more weakly (19 per cent and $+$0.08).

These patterns should be borne in mind when interpreting findings, particularly for the
visit-based tier. Warmth and strictness estimates for the primary tier
(\Cref{sec:p1_data}) may not be fully generalisable to schools in the North of England,
and cross-tier comparisons should account for the possibility that regional compositional
differences contribute to attenuation.


\begin{table}[htbp]
\centering
\small
\caption{Comparison of interviewed and visited schools to the national population of open secondary schools in England.}
\label{tab:2A:representativeness}
\begin{tabular}{ll ccc cc}
\toprule
 & & \multicolumn{3}{c}{\textit{Percentage of schools} (\%)} & \multicolumn{2}{c}{\textit{p}-value} \\
\cmidrule(lr){3-5} \cmidrule(lr){6-7}
\textit{Variable} & \textit{Category} & National & Interviewed & Visited & Int. & Vis. \\
 & & $(n=3{,}312)$ & $(n=300)$ & $(n=103)$ & & \\
\midrule
  \textit{School type} & Academy & 74.7 & 72.3 & 74.3 & \textit{ns} & \textit{ns} \\
   & Free school & 9.2 & 10.0 & 11.9 &  &  \\
   & Maintained & 15.8 & 17.7 & 13.9 &  &  \\
\addlinespace[3pt]
  \textit{Phase} & Secondary & 95.1 & 94.3 & 93.1 & \textit{ns} & \textit{ns} \\
   & All-through & 4.9 & 5.7 & 6.9 &  &  \\
\addlinespace[3pt]
  \textit{Gender} & Mixed & 89.2 & 87.7 & 85.1 & \textit{ns} & \textit{ns} \\
   & Girls & 6.1 & 7.0 & 5.9 &  &  \\
   & Boys & 4.3 & 5.3 & 8.9 &  &  \\
\addlinespace[3pt]
  \textit{Admissions policy} & Non-selective & 85.4 & 87.0 & 88.1 & \textit{ns} & \textit{ns} \\
   & Selective & 4.9 & 5.3 & 5.0 &  &  \\
\addlinespace[3pt]
  \textit{Urban/rural} & Urban & 87.4 & 88.3 & 90.1 & \textit{ns} & \textit{ns} \\
   & Rural & 12.5 & 11.7 & 9.9 &  &  \\
\addlinespace[3pt]
  \textit{Has sixth form} & Yes & 62.7 & 68.0 & 75.2 & \textit{ns} & p$<$.01 \\
   & No & 37.3 & 32.0 & 24.8 &  &  \\
\addlinespace[3pt]
  \textit{Ofsted grade (Aug 2024)} & Outstanding & 14.7 & 18.3 & 25.2 & p$<$.05 & p$<$.05 \\
   & Good & 69.2 & 70.3 & 63.1 &  &  \\
   & Requires improvement & 12.9 & 8.0 & 8.7 &  &  \\
   & Inadequate & 3.2 & 3.3 & 2.9 &  &  \\
\addlinespace[4pt]
\multicolumn{7}{l}{\textit{Continuous variables}} \\
\addlinespace[2pt]
  \textit{Number of pupils} & Mean (\textit{SD}) & 1{,}079 (396) & 1{,}109 (384) & 1{,}134 (380) & \textit{ns} & \textit{ns} \\
  \textit{FSM percentage} & Mean (\textit{SD}) & 29.9 (14.7) & 28.0 (14.7) & 29.1 (15.7) & p$<$.05 & \textit{ns} \\
  \textit{Prior P8 (2018--19)} & Mean & $+$0.04 & $+$0.08 & $+$0.17 & \textit{ns} & p$<$.01 \\
\bottomrule
\end{tabular}
\begin{minipage}{\linewidth}
\smallskip
\footnotesize\textit{Notes:} National population is all open secondary and all-through schools in England.
$p$-values are approximate: $\chi^2$ goodness-of-fit test (observed vs national proportions) for categorical
variables; one-sample $z$-test for continuous variables. Rows with $<$0.5\% share in all columns omitted.
The Ofsted-grade and prior-P8 rows are computed on the full tiers (visited $n=103$, all
graded; interviewed $n=303$, of which 300 graded; national $n=3{,}332$, of which 3,297
graded); grade tests are $\chi^2$ against the non-tier remainder (visited $p=0.019$,
interviewed $p=0.029$), and the prior-P8 rows use schools with a published 2018--19
score (national $n=2{,}731$, interviewed $n=259$, visited $n=95$; two-sample $t$-tests,
visited $p=0.003$, interviewed $p=0.087$).
Regional distribution is shown separately in \cref{fig:2A:representativeness_region}.
\end{minipage}
\end{table}

\begin{figure}[htbp]
  \centering
  \caption{Regional distribution of interviewed and visited schools}
  \label{fig:2A:representativeness_region}
  \includegraphics[width=\textwidth]{fig_D_regional_representativeness.pdf}
\begin{minipage}{\linewidth}\smallskip
\footnotesize\textit{Notes:} The comparison is with the national population of open secondary and all-through schools in England. Bars show the percentage of schools in each region for the national population ($n = 3{,}312$), interviewed schools ($n = 300$), and visited schools ($n = 102$), the schools matched to the school panel (the 2023/24 regional characteristics panel holds 3{,}312 of the 3{,}332 analysis-panel schools; the remainder lack a region record at that date). Regions are ordered by national share.
\end{minipage}
\end{figure}

\paragraph{Coverage of the text collections.} The three collections of
text were each gathered from the register of open state secondary schools at
a different date, which is why their raw counts differ. Inspection reports
were scored for 3,324 schools, behaviour policies were obtained for 3,363,
and websites were crawled for 3,370. Set against the analysis panel of 3,332 schools, the coverage counts and shares are those reported with the data: 3,272 reports, 3,294 policies and 3,311 crawls, of which 3,284 contain scorable prose. The few schools in each collection that
are outside the panel are schools that closed, merged or changed identifier
between the crawl and the construction of the panel.

\clearpage
% - - - - - - - - - - - - - - - - - - - - - - - - - - - - - - -
\section{Prediction Model Specification}
\label{sec:p1_app_prediction}

The text representation is TF-IDF over words and two-word phrases
(English stop words removed; terms kept if they appear in at least five
documents and at most ninety per cent of them; sublinear term frequency;
each document's vector normalised to unit length). The model is a ridge
regression with an unpenalised intercept, estimated exactly in kernel
form. Every reported correlation is leave-one-out: the model is re-estimated with each school excluded, and the penalty is re-chosen within every fold
by closed-form inner leave-one-out, so the held-out school takes no part in estimating the model or choosing its penalty. Predictions are clipped to the range of the
training scores. Significance comes from a permutation test: the identical
pipeline is run on two hundred permutations of the outcome, and the
reported $p$ is the share of permutations matching or beating the real
correlation.

Before modelling, documents are cleaned of identifying and clerical
material: establishment, trust and local-authority names (from the
national register), month names and years, postcodes, tokens containing
digits, inspection-report page furniture, letter salutations, and
framework boilerplate. The demographic adjustment regresses each outcome on the school's free-school-meals share, share of pupils with English as an additional language, prior attainment and report age, and takes the residual as the target. The mark-scheme bands are not changed. The adjustment is applied to the outcome only. The adjusted figure for a mark scheme is the correlation between the band as awarded and this residual.

Alternative representations were evaluated in the identical harness and none improved on the specification above. In brief, adding three- and four-word phrases moved each correlation by less than its sampling variation. Topic-model features predicted every dimension worse than word and phrase counts did. Models estimated on passage embeddings from two pretrained encoders correlate with enacted warmth and teaching at essentially zero. They do reproduce the espoused-strictness result. So the warmth and teaching signals lie in the particular words used, not in meaning that survives paraphrase. One result disagrees between the two encoders. The larger encoder predicts enacted strictness at $+0.33$ with mean pooling. The smaller encoder's correlation is near zero, and so is that of every word-based model. Because it appears under one encoder and not the other, it is recorded here as unresolved rather than used.
Models built on sentence tags alone (topic and polarity) predicted every score worse than models built on the words. A second variant multiplied each word count by plus one in a positive sentence and minus one in a negative sentence. This raised the correlation for warmth, strictness and teaching alike. The words responsible were clerical ones from the report format. A feature that helps all three dimensions at once, and does so through format words, is picking up how the report was written rather than what the school is like. The chapter's rule discards such features, and it was discarded.

\section{What the Models Use: Examples}
\label{sec:p1_app_examples}

This appendix supports the validity of the warmth and teaching prediction
models by showing their evidence in readable units. The models work on
weighted word and phrase counts, but a sentence's score is simply the sum
of its phrases' weights, so every sentence in the corpus can be ranked by
how strongly it pushes a school's score up or down. The listings below
give the strongest sentences in each direction for each model, followed by the strongest longer phrases (\cref{tab:p1_phrases}). The pattern is consistent: sentences
describing adults attending to pupils, and teaching designed and checked,
raise the scores; stock assurances and provision boilerplate lower them.
One oddity in the teaching listing should be read with care: several of
the score-raising sentences are criticisms (``pupils are not remembering
knowledge deeply''). A report mentions long-term memory, whether in praise or criticism, only when the school teaches in that way. The model counts words and does not register whether a sentence affirms or denies. For teaching this does little harm, because the vocabulary itself is informative. For strictness it is fatal. The informative strictness sentences are the ones that say behaviour is not good, and the model cannot tell those from the ones that say it is.

\subsection*{Warmth}
\paragraph{Sentences that raise the score.}
\begin{itemize}
  \item ``They are taught how to show respect for different opinions and listen to the views of others.''
  \item ``They hold leaders strictly to account for the quality of the school's provision and are highly effective.''
  \item ``They listen to others and express their views in a well-managed environment.''
  \item ``Sixth-form students are proud role models and make a highly positive contribution to the school and local community.''
  \item ``Those responsible for governance know the school's context and provide support and professional challenge where necessary to ensure that the school's high standards are maintained.''
  \item ``In lessons, pupils are diligent, focused, and eager to learn.''
  \item ``They listen to staff and take into account the pressures on them.''
  \item ``Pupils are encouraged to contribute to the school community.''
  \item ``In lessons, pupils listen respectfully to teachers and each other.''
  \item ``The school's motto is `excellence with care.' Pupils testify to how this works in practice.''
  \item ``In the secondary phase, pupils receive regular careers information and guidance.''
  \item ``They move calmly between classes and wear their uniforms smartly, meeting the school's high standards.''
\end{itemize}
\paragraph{Sentences that lower the score.}
\begin{itemize}
  \item ``However, some pupils, particularly some disadvantaged pupils, are absent too often.''
  \item ``Many pupils feel confident they have a `trusted adult' they can go to.''
  \item ``(Information for the school and appropriate authority)  In some subjects, the work given to some pupils does not enable them to develop their ideas and extend their answers.''
  \item ``Pupils take part in a range of extra-curricular activities to develop their talents and interests, for example basketball, choir, photography, and cookery.''
  \item `` Sometimes, the work given to pupils, including those with SEND, does not enable them to achieve the aims and ambition of the curriculum.''
  \item ``However, in some subjects, the work given to some pupils does not enable them to develop their ideas and extend their answers as well as they could.''
  \item ``Missing out on important learning means that these pupils do not achieve as well as they might.''
  \item ``Pastoral staff work closely with these pupils and their families.''
  \item ``They check closely to ensure that disadvantaged pupils have the support they need.''
  \item ``The work given to pupils does not always give them opportunities to develop deeper thinking, such as reasoning and justifying.''
  \item ``Leaders monitor closely pupils' participation in the school's extra-curricular offer.''
  \item ``The school recognises the importance of ensuring that pupils, particularly disadvantaged pupils, and pupils with SEND, gain wider opportunities beyond the curriculum.''
\end{itemize}
\paragraph{The strongest longer phrases.} Raising: ``high expectations behaviour''; ``early career teachers''; ``provide effective support''; ``pupils attend school regularly''; ``school high expectations''; ``positive attitudes learning''; ``meet pupils needs''; ``pupils need help''. Lowering: ``work given pupils''; ``support pupils learning''; ``develop leadership skills''; ``expectations pupils achieve''; ``inspection headteacher school''; ``pupils kept safe''; ``broad ambitious curriculum''; ``ambitious curriculum pupils''.
\subsection*{Teaching}
\paragraph{Sentences that raise the score.}
\begin{itemize}
  \item ``Those responsible for governance know the school's context and provide support and professional challenge where necessary to ensure that the school's high standards are maintained.''
  \item ``Leaders are supported and challenged by dedicated and knowledgeable governors, who share their vision of high standards in a supportive school.''
  \item ``In a few subjects, however, pupils are not remembering knowledge deeply in the long term.''
  \item ``The school's motto is `excellence with care.' Pupils testify to how this works in practice.''
  \item ``They lead charity initiatives and community-based projects to instil a keen sense of responsibility and empathy for others.''
  \item ``In lessons pupils settle to their work quickly and show high levels of respect to their teachers and to other pupils.''
  \item ``They are taught how to show respect for different opinions and listen to the views of others.''
  \item ``In lessons, most pupils focus on their learning and learn well.''
  \item ``They carefully choose the core knowledge they want pupils to know and remember in the long term.''
  \item ``Staff, pupils, and parents share a strong sense of pride in belonging to this school's community.''
  \item ``Leaders have ensured that all aspects of the national curriculum are taught in Years to Work with primary schools helps new pupils to settle in quickly.''
  \item ``Teachers help pupils to embed learning in their long-term memory.''
\end{itemize}
\paragraph{Sentences that lower the score.}
\begin{itemize}
  \item ``The personal, social and health education (PSHE) curriculum is designed to support pupils' wider development.''
  \item ``Pupils take part in a range of extra-curricular activities to develop their talents and interests, for example basketball, choir, photography, and cookery.''
  \item ``On other occasions, adults do not support pupils to complete work to a high standard.''
  \item ``Judgements in this report are based on the current inspection framework and also reflect changes that have happened at any point since the last graded inspection.''
  \item ``The school offers pupils a range of ways to develop their wider skills and talents.''
  \item ``The school has carefully designed the personal, social and health education curriculum so that it supports pupils' wider development.''
  \item ``It does so, not only by having high expectations for pupils, but by supporting their individual aspirations and talents.''
  \item ``The curriculum in personal, social and health education is designed to help pupils understand important ideas.''
  \item ``Well-planned personal, social, health education ensures that pupils learn about relevant and sensitive issues, including personal safety, relationships, sexual health and consent.''
  \item ``However, some pupils, particularly some disadvantaged pupils, are absent too often.''
  \item ``They have an accurate understanding of what remains to be done to achieve this for all pupils.''
  \item ``Pastoral staff work closely with these pupils and their families.''
\end{itemize}
\paragraph{The strongest longer phrases.} Raising: ``early career teachers''; ``school high expectations''; ``high expectations behaviour''; ``important knowledge pupils''; ``behave lessons school''; ``pupils need additional''; ``positive attitudes learning''; ``key knowledge skills''. Lowering: ``support pupils learning''; ``develop leadership skills''; ``school community pupils''; ``pupils kept safe''; ``high expectations pupils achieve''; ``ambitious curriculum pupils''; ``expectations pupils achieve''; ``school pupils enjoy''.


% generated by make_ch2_exhibits.py -- do not edit
\begin{table}[htbp]
\centering
\caption{The phrases with the largest weights}
\label{tab:p1_phrases}
{\small\begin{tabular}{llll}
\toprule
\multicolumn{2}{c}{Warmth} & \multicolumn{2}{c}{Teaching} \\
\cmidrule(lr){1-2}\cmidrule(lr){3-4}
Raises & Lowers & Raises & Lowers \\
\midrule
listen & closely & writing & standard \\
lessons pupils & welcoming & lessons pupils & case \\
small number & standard & listen & staff work \\
order & feel confident & share & health \\
exemplary & given pupils & standards & expectations pupils \\
lesson & pupils particularly & explain & pupils achieve \\
guidance & expectations pupils & highly & curriculum school \\
leaders check & including send & lesson & disadvantaged pupils \\
managed & international & motto & pupils kept \\
contribute & case & students benefit & broad ambitious \\
\bottomrule
\end{tabular}}
\begin{minipage}{\linewidth}
\smallskip
\footnotesize\textit{Notes:} What the prediction models use: the ten phrases with the largest positive and negative weights in the warmth and teaching models (inspection reports, demographics removed from the target). Full-sample descriptive estimates at the penalty the tested models chose.
\end{minipage}
\end{table}


\clearpage
\section{Further Parent-Comparison Checks}
\label{sec:p1_app_parents}

This section collects the checks that qualify the parent comparison of \cref{sec:p1_pv}: what the visited schools show, a stricter form of the grade check, and a test that the report text does not contain the parents' answers in a form the model uses.

The visited schools complicate the picture. There the parent measure can be checked against both
what the observers saw and what the head said. Against enacted warmth it correlates at $+0.17$, positive but too imprecise to rely on ($p = 0.08$, $n = 101$).
Against the head's espoused warmth it correlates at $+0.43$. Parents agree with the head more than
with the observers. There are two ways to interpret that and the data cannot tell us which is
correct.
The first is that parents and heads get a picture of the school across an entire year or longer,
while the observers saw just one day. On that basis the parents' views provide evidence for the head's account. The second possibility is that the parents and the head, both invested in the school, may report similar things, whatever the school is actually like in practice. The survey therefore cannot
judge between what schools say and what observers see.

The check in the main text uses the grade from an earlier inspection. A stricter version uses the grade awarded by the same inspection whose report the model reads. Comparing schools with the same such grade halves the warmth correlation, to $+0.16$. The grade and the prose come from one visit, so the grade already reflects some of the warmth the prose describes. Removing it therefore removes some warmth as well as a rival explanation. The
national parent correlation also falls by about half when the grade is held constant,
from $+0.29$ to $+0.14$. For scale, the grade alone says little about warmth: on the
visited schools it correlates with enacted warmth at only $-0.20$ (grades run from 1,
Outstanding, to 4, so a negative sign means better-graded schools are warmer). The grade alone is only weakly related to warmth. The text prediction is therefore using information beyond the grade.

One caveat needed testing. Inspectors read Parent View during inspections, and about 3 per cent
of reports quote it. The correlation with parents could therefore arise because the report text contains the parents' own answers. Deleting every sentence that mentions the survey and re-running the whole prediction
leaves the correlation unchanged at $+0.29$. The quoted survey material is not what drives the result.

\section{The Mark Schemes}
\label{sec:p1_app_rubrics}

The seven schemes below are reproduced word for word. They were finalised before any document was scored and not changed afterwards.
Each was applied once per document by gpt-4o-mini (version-pinned), which
returned a band and one quoted evidence sentence. Thirty documents per source were re-scored with caching disabled, with the agreement figures reported in the chapter, and the full grid was re-scored
once by gpt-4o as a sensitivity check, reported alongside and never used
to select results.

\noindent The pinned model versions: every scheme's primary scorer is
\texttt{gpt-4o-mini}, pinned to the dated snapshot
\texttt{gpt-4o-mini-2024-07-18}; the whole-grid sensitivity re-score uses
\texttt{gpt-4o}, pinned to \texttt{gpt-4o-2024-08-06}
(\texttt{model\_pins.py}, verified against the provider's model list before
scoring).\medskip

\subsection*{1. Inspection reports --- Warmth}
What is measured: how the report describes the relationships between staff and pupils. Score the relational quality described, never the pastoral provision listed.

\begin{enumerate}
  \item Relationships described as poor: pupils unknown to staff, uncared for, or unsafe because of how staff relate to them.
  \item Any credible negative relational signal --- pupils not always heard, bullying unresolved, staff not accessible --- even alongside positives.
  \item DEFAULT. A safe and generally positive school, described through labels, provision or outcomes ("caring school", "pupils feel safe", "pastoral team is effective") without describing how staff and pupils actually relate.
  \item The report describes the relationships themselves: staff know pupils well, warm interactions observed, pupils feel valued by staff as individuals.
  \item Relational warmth described as a defining feature of the school, in emphatic terms, across several passages.
\end{enumerate}
Rules: if the report does not discuss relationships, score 3 --- silence is not coldness. The 3/4 line: does the sentence describe how staff and pupils relate, or the school's character and what provision exists? In doubt between two bands, take the lower.

\subsection*{2. Inspection reports --- Strictness}
What is measured: what the report says about conduct standards and how consistently the adults enforce them.

\begin{enumerate}
  \item Disruption is a persistent, documented problem; rules unclear or not applied.
  \item Standards exist but enforcement is patchy; behaviour is named as an area for improvement, or inspectors saw low-level disruption.
  \item DEFAULT. Clear expectations, generally maintained; behaviour mentioned but unremarkable.
  \item The report as a whole presents conduct as a strength, in lessons and around the building.
  \item Conduct is exceptional and unmistakably a defining feature of the school.
\end{enumerate}
Rules: if behaviour is simply not discussed, score 3 --- a 1 or 2 needs positive evidence of a problem. Pupils' opinions about the rules are warmth evidence, never strictness evidence: strictness is what the adults do. In doubt between two bands, take the lower.

\subsection*{3. Inspection reports --- Teaching}
What is measured: how the report describes classroom teaching --- explanations, subject knowledge, and the checking of understanding. Curriculum-design language (intent, sequencing) is background, not evidence.

\begin{enumerate}
  \item Teaching named as a serious failure and a cause of poor outcomes.
  \item Teaching inconsistent ("some teachers... not always") and named as requiring improvement.
  \item DEFAULT. Competent and unremarkable, or classroom delivery not discussed at all.
  \item Teaching described positively and with specifics across most of the school ("typically", "most teachers").
  \item Teaching described as expert and a defining strength ("consistently", "across all subjects").
\end{enumerate}
Rules: score how teachers teach, not what the curriculum intends. The consistency words carry the ladder: "consistently" points to 5, "typically" to 4, "some" to 2. In doubt, or with no delivery evidence, score 3.

\subsection*{4. Behaviour policies --- Strictness}
What is measured: how far the policy removes staff discretion and prescribes pupil conduct. Score the machinery the policy describes, not how severe its language sounds.

\begin{enumerate}
  \item No consequence language at all --- the policy is aspirational or entirely about rewards.
  \item Consequence language is present, but no day-to-day system a new member of staff could act on tomorrow.
  \item DEFAULT. A working day-to-day system, applied by staff judgement from a menu of options. The standard English secondary policy.
  \item A named tracking or tier system: a pupil's position in the system, not staff judgement, decides which consequence follows.
  \item As 4, and the policy prescribes a default deportment norm --- a SLANT-type routine or a similar required way of sitting, attending or moving.
\end{enumerate}
Rules: the test for 3 --- could a new teacher read this policy and know what to do tomorrow when a lesson is disrupted? Harsh penalties at staff discretion are LESS strict on this construct than mild ones applied automatically.

\subsection*{5. Behaviour policies --- Warmth}
What is measured: what the policy tells staff about how to relate to pupils. Rewards and pastoral structures are provision, not warmth.

\begin{enumerate}
  \item No positive or relational content beyond the sanctions system.
  \item Positive values asserted ("every child matters") with no described mechanism and no guidance to staff.
  \item DEFAULT. A positive mechanism described --- named rewards with the occasions they are given on, or named pastoral routes --- but no guidance to staff on how to relate to pupils. Most policies.
  \item Explicit operational guidance to staff on relating to pupils as individuals: greet at the door, learn names, how to hold a repair conversation.
  \item As 4, and relational quality is the organising principle: consequences framed as relationship-preserving, and the policy says the relationship survives the sanction.
\end{enumerate}
Rules: a scheduled restorative meeting is machinery; only telling staff how to conduct it is guidance. If in doubt whether a sentence is guidance, it is not --- a rich reward economy cannot pass 3.

\subsection*{6. Websites --- Strictness}
What is measured: whether the site's high-expectations language is about behaviour at all, and how far the school goes in claiming strictness of itself.

\begin{enumerate}
  \item No mention of high standards or high expectations anywhere --- even a cursory mention moves it up.
  \item High standards or expectations claimed, but generic --- not tied to behaviour or conduct.
  \item DEFAULT. High expectations explicitly about behaviour or conduct --- or the school asserting it is demanding, making no apologies.
  \item Behaviour treated as central: expected at all times, pupils required to become self-disciplined, rewards-and-sanctions machinery on the site, or behaviour named as the foundation of success.
  \item The school states that it is strict: strict routines and protocols, zero tolerance, no excuses, the highest possible standards.
\end{enumerate}
Rules: aspiration language alone ("achieve their full potential") does not reach band 2. It is the explicit statement of strictness that earns the 5, not the severity of the machinery.

\subsection*{7. Websites --- Warmth}
What is measured: how far the site presents care for pupils and staff--pupil relationships as what the school is about, in the school's own words. The axis: absent --- courtesy --- claimed --- stressed --- central.

\begin{enumerate}
  \item No care or support language at all. The school presents itself entirely through results, facilities, opportunities and ambition.
  \item The courtesy register only: generic supportive-atmosphere phrases ("a supportive learning environment", "care and support for each individual") appearing in passing, typically once inside the head's welcome, and not returned to.
  \item DEFAULT. Care claimed as a feature of the school: pastoral care named, support for individual pupils emphasised, or a caring atmosphere claimed prominently or in more than one place.
  \item Care or relationships stressed as important and returned to: pupils described as known and valued as individuals, care presented as part of what the school is for, across more than one page or passage.
  \item Staff--pupil relationships explicitly named and presented as central to what the school is --- and more than once. Rare by design.
\end{enumerate}
Rules: the 2/3 line is courtesy versus feature --- if the care language lives only in the head's welcome paragraph, score 2. Band 5 needs the relationships themselves named (not just "caring" or "supportive"), plus centrality, plus recurrence. Quoted Ofsted praise does not count; the school's own and its trust's words do. In doubt between two bands, take the lower.


